% =====================================================================
%  figset_B.tex  --  TikZ figure set B of the QCD factorization monograph
%  (agent figB, local style prefix "fb").
%  Scope (FIGURES_INDEX_v2):
%    Sec. 2 -- Ch9 method, Ch10 disall, Ch11 coll (required + optional);
%    Sec. 3 -- Ch4 ir, Ch6 ee, Ch7 dis, Ch9 method, Ch10 disall
%              (required + optional).
%  Figures, in book order (owner chapter in brackets; "opt" = optional):
%    [Ch4]  fig:ir-mode-exponents
%           fig:ir-landau-simplex          (opt; Sec. 2 optional table, Ch4
%                                           owner -- drawn because no other
%                                           set covers Ch4; may be dropped)
%    [Ch6]  fig:ee-cutsum-disc
%    [Ch7]  fig:dis-euclid-segment
%           (the optional fig:dis-twist-graphical lives in figset_F.tex; the
%           duplicate drawn here was removed by the figure checker)
%    [Ch9]  fig:method-flowchart, fig:method-hasse, fig:method-nesting,
%           fig:method-neighbourhoods, fig:method-kg, fig:method-sudakov-plane,
%           fig:method-coeff-minus-coll, fig:method-glauber-direction
%    [Ch10] fig:disall-seed, fig:disall-ladder-regions, fig:disall-heavy
%    [Ch11] fig:coll-sia, fig:coll-sidis, fig:coll-photoproduction,
%           fig:coll-transversity (opt)
%
%  Requirements: common/preamble_fact.tex (ctexbook, XeLaTeX).  It loads
%  tikz with the libraries decorations.pathmorphing, decorations.markings,
%  arrows.meta, calc, positioning, shapes.geometric, patterns, backgrounds,
%  fit, and pgfplots (compat=1.17), and contains the style blocks STYLES and
%  STYLES-EXT (quark, gluon, gluon loop, photon, ghost, mom, mom', vertex,
%  blob, wilson, wilson plain, cut, hard blob, jet blob, jetB blob, soft blob,
%  hadron blob, glauber, opvertex, scalar, longit + \longitmark, lcaxis,
%  region), the macros \ii \dd \MSbar \as \order \Real \Imag \Disc, the
%  xcolor option dvipsnames (OliveGreen, RoyalBlue, Plum) and cleveref.
%  No package or TikZ library is loaded here.  Every extra style is declared
%  in the options of its tikzpicture (local to that picture) and carries the
%  prefix "fb"; the few helper macros (\fbEEladder, \fbEEred, \fbref,
%  \fbstep, \fbKGframe, \fbCMCphot, \fbCMCapprox, \fbGwedge, \fbTWframe)
%  are defined inside their figure environment, hence local to it.
%
%  Conventions used throughout (WRITER_SPEC_v2 Sec. 3):
%    * solid line with arrow = quark; the arrow is the fermion-number flow
%      (for an antiquark it points against the momentum); thick solid line
%      with arrow = heavy quark (fig:disall-heavy);
%    * separate thin labelled arrows (mom/mom') = momentum flow;
%    * double line with arrow = Wilson line, the arrow is the path
%      orientation of W[b,a] (from a to b); future pointing = away from the
%      operator vertex towards infinity;
%    * red dashed line = final-state cut (for a vertical cut the amplitude
%      is on the left, the complex-conjugate amplitude on the right);
%      zig-zag = branch cut;
%    * blob colours as in STYLES-EXT: grey = hard, blue = collinear to A
%      (target), orange = collinear to B / current jet, green = soft,
%      light grey = hadron; dark-red coil = Glauber gluon; blue coil with a
%      dot = scalar-polarized (K) gluon, the dot marking its k^mu end;
%    * n_A, n_B = light-like directions of the two beams/jets (row 2),
%      e_G = n_B - n_A (row 38); light-cone and exponent axes are labelled.
%  Each figure is delimited by  %%% BEGIN <label>  /  %%% END <label>.
% =====================================================================


%%% BEGIN fig:ir-mode-exponents
\begin{figure}[htbp]
\centering
\begin{tikzpicture}[fd,x={(1.9cm,0cm)},y={(-0.996cm,0.575cm)},z={(0cm,1.55cm)},
  fb ax/.style={lcaxis,gray!50!black},
  fb guide/.style={draw=gray!60,line width=0.4pt,dash pattern=on 1pt off 1.6pt},
  fb plane/.style={fill=gray!22,fill opacity=0.62,draw=gray!55!black,line width=0.6pt},
  fb pt/.style={circle,draw=white,line width=0.4pt,inner sep=0pt,minimum size=5.6pt},
  fb gl/.style={circle,draw=red!60!black,fill=white,line width=0.9pt,inner sep=0pt,minimum size=6pt},
  fb lab/.style={font=\footnotesize,inner sep=2pt},
  fb tick/.style={font=\scriptsize,text=gray!35!black,inner sep=1.5pt}]
  % ---------------- base grid (c = 0) and axes ----------------
  \foreach \t in {1,2} {
    \draw[fb guide] (\t,0,0) -- (\t,2,0);
    \draw[fb guide] (0,\t,0) -- (2,\t,0);
  }
  \draw[fb ax] (0,0,0) -- (2.75,0,0) node[right,font=\small] {$a$};
  \draw[fb ax] (0,0,0) -- (0,2.85,0) node[left,font=\small] {$b$};
  \foreach \t in {1,2} {
    \draw[gray!50!black,line width=0.5pt] (\t,0,0) -- ++(0,0,-0.07) node[below,fb tick] {$\t$};
    \draw[gray!50!black,line width=0.5pt] (0,\t,0) -- ++(0,0,-0.07) node[below,fb tick] {$\t$};
  }
  % vertical guides down to the base plane
  \draw[fb guide] (2,2,0) -- (2,2,0.95);
  \draw[fb guide] (2,0,0) -- (2,0,1);
  \draw[fb guide] (0,2,0) -- (0,2,1);
  % Glauber point and its vertical offset: below the plane, hence drawn before the translucent plane
  \node[fb gl] (G) at (2,2,1) {};
  \draw[red!60!black,line width=0.85pt,dash pattern=on 2.5pt off 1.8pt,-{Stealth[length=4pt,width=3.4pt]}]
       (2,2,1.07) -- (2,2,1.93);
  % ---------------- on-shell plane c = (a+b)/2 ----------------
  \draw[fb plane] (0,0,0) -- (2,0,1) -- (2,2,2) -- (0,2,1) -- cycle;
  \node[font=\footnotesize,text=gray!30!black,rotate=22.2] at (1.35,0.3,0.825) {在壳平面};
  % the c axis lies in front of the plane
  \draw[fb ax] (0,0,0) -- (0,0,2.8) node[above,font=\small] {$c$};
  \foreach \t in {1,2}
    \draw[gray!50!black,line width=0.5pt] (0,0,\t) -- ++(-0.045,0,0) node[left,fb tick] {$\t$};
  % ---------------- modes ----------------
  \node[fb pt,fill=black] (H)  at (0,0,0) {};
  \node[fb pt,fill=blue!55!black] (CA) at (0,2,1) {};
  \node[fb pt,fill=orange!85!black] (CB) at (2,0,1) {};
  \node[fb pt,fill=green!40!black] (US) at (2,2,2) {};
  \node[fb pt,fill=green!60!black] (S) at (1,1,1) {};
  \node[fb lab,anchor=north east] at (H.south west) {硬 $(0,0,0)$};
  \node[fb lab,anchor=east] at (CA.west) {$A$ 共线 $(0,2,1)$};
  \node[fb lab,anchor=west] at (CB.east) {$B$ 共线 $(2,0,1)$};
  \node[fb lab,anchor=west] at (US.east) {超软 $(2,2,2)$};
  \node[fb lab,anchor=south] at ($(S.north)+(0cm,0.02cm)$) {软 $(1,1,1)$};
  \node[fb lab,anchor=west,text=red!55!black] (GL) at (canvas cs:x=3.15cm,y=3.25cm) {Glauber $(2,2,1)$};
  \node[fb lab,anchor=north west,text=red!55!black,inner sep=1pt] at ($(GL.south west)+(0.1cm,0cm)$) {$k^2\approx-k_T^2$，离壳};
  \draw[red!55!black,line width=0.4pt] ($(G.east)+(0.02cm,0.03cm)$) -- (GL.west);
  % ---------------- legend ----------------
  \node[font=\small,anchor=east] at (canvas cs:x=-0.45cm,y=4.3cm) {$(k^+,\,k^-,\,k_T)\sim(\lambda^a,\,\lambda^b,\,\lambda^c)\,Q$};
  \node[font=\footnotesize,anchor=north west,align=left,inner sep=0pt] at (canvas cs:x=6.45cm,y=3.38cm)
    {平面上 $c=\tfrac12(a+b)$：\\[1pt]$\quad k_T^2\sim k^+k^-\sim|k^2|$\\[5pt]
     平面下 $c<\tfrac12(a+b)$：\\[1pt]$\quad k_T^2\gg|k^+k^-|$，$k^2\approx-k_T^2$};
\end{tikzpicture}
\caption{动量模式在指数空间 $(a,b,c)$ 中的位置，$(k^+,k^-,k_T)\sim(\lambda^a,\lambda^b,\lambda^c)\,Q$，$\lambda\ll1$（指数越大，相应分量越小）。
硬、$A$ 共线、$B$ 共线、软与超软模式都在在壳平面 $c=\frac12(a+b)$（半透明灰色平行四边形）上，即 $k_T^2$、$k^+k^-$ 与虚度 $|k^2|$ 同阶；
硬、两种共线与超软模式恰为平行四边形的四个顶点，软模式 $(1,1,1)$（TMD 因子化所用的标度）为其中心。
Glauber 模式 $(2,2,1)$（红色空心圆，透过平面看到）与超软模式有相同的 $k^\pm$，只是横动量大 $\lambda^{-1}$ 倍（红色虚线），因而位于平面之下：
$|k^+k^-|\ll k_T^2$，$k^2\approx-k_T^2$，以横动量为主且离壳；在 $(\ln k^+,\ln k^-)$ 平面（\cref{fig:momentum-regions}）上它与超软模式重合，只有加上第三个方向 $\ln k_T$ 才能区分。
本书的“超软”即 Libby--Sterman 的“软”。}
\label{fig:ir-mode-exponents}
\end{figure}
%%% END fig:ir-mode-exponents


%%% BEGIN fig:ir-landau-simplex
\begin{figure}[htbp]
\centering
\begin{tikzpicture}[fd,
  fb lab/.style={font=\footnotesize,inner sep=1.5pt},
  fb edgeA/.style={draw=red!65!black,line width=2.2pt,line cap=round},
  fb edgeB/.style={draw=blue!65!black,line width=2.2pt,line cap=round},
  fb simplex/.style={fill=gray!10,draw=gray!60!black,line width=0.6pt},
  fb corner/.style={circle,fill=black,inner sep=0pt,minimum size=5.5pt}]
  % ======================= the vertex graph and its lines =======================
  \begin{scope}[yshift=1.25cm]
    \coordinate (X) at (0,0); \coordinate (Y1) at (1.3,0.8); \coordinate (Y2) at (1.3,-0.8);
    \draw[photon] (-1.2,0) -- (X);
    \draw[quark] (X) -- (Y1); \draw[quark] (Y1) -- (2.15,1.32);
    \draw[quark] (2.15,-1.32) -- (Y2); \draw[quark] (Y2) -- (X);
    \draw[gluon] (Y2) -- (Y1);
    \foreach \v in {X,Y1,Y2} \node[vertex] at (\v) {};
    \node[fb lab,anchor=south] at (-0.62,0.08) {$q$};
    \node[fb lab,anchor=south east] at (0.66,0.4) {$p_1-k$};
    \node[fb lab,anchor=north east] at (0.66,-0.4) {$p_2+k$};
    \node[fb lab,anchor=west] at (1.42,0) {$k$};
    \node[fb lab,anchor=west] at (2.15,1.32) {$p_1$};
    \node[fb lab,anchor=west] at (2.15,-1.32) {$p_2$};
    \node[font=\footnotesize,align=center,anchor=north] at (0.5,-1.72) {线 $0,1,2$：$k$，$p_1-k$，$p_2+k$\\Feynman 参数 $\alpha_0,\alpha_1,\alpha_2$};
  \end{scope}
  % ======================= (a) on-shell massless vertex =======================
  \begin{scope}[xshift=4.85cm,x=2.55cm,y=2.55cm]
    \coordinate (A0) at (0,0); \coordinate (A1) at (1,0); \coordinate (A2) at (0.5,0.866);
    \draw[fb simplex] (A0) -- (A1) -- (A2) -- cycle;
    \draw[gray!55!black,line width=0.9pt,dash pattern=on 3pt off 2pt] (A1) -- (A2);
    \draw[fb edgeA] (A0) -- (A1);
    \draw[fb edgeB] (A0) -- (A2);
    \node[fb corner] at (A0) {};
    \node[fb lab,anchor=north east] at (A0) {$S$\,($\alpha_0=1$)};
    \node[fb lab,anchor=north west] at (A1) {$\alpha_1=1$};
    \node[fb lab,anchor=south] at (A2) {$\alpha_2=1$};
    \node[fb lab,anchor=north,text=red!60!black] at (0.5,-0.03) {$C_1$};
    \node[fb lab,anchor=south east,text=blue!60!black] at (0.23,0.42) {$C_2$};
    \node[fb lab,anchor=south west,text=gray!35!black] at (0.77,0.42) {$Q^2=0$ 时};
    \node[font=\small] at (0.5,0.3) {$H$};
    \node[fb lab,anchor=north,align=left] at (0.5,-0.2) {$C_1$：$\alpha_2=0$，$k=xp_1$\\$C_2$：$\alpha_1=0$，$k=-yp_2$\\$S$：$k=0$；\ $H$：没有钳夹};
    \node[font=\small,anchor=north] at (0.5,-0.75) {(a)\ 在壳无质量：$\mathcal F=-Q^2\alpha_1\alpha_2$};
  \end{scope}
  % ======================= (b) off-shell (Euclidean) triangle =======================
  \begin{scope}[xshift=10.55cm,x=2.55cm,y=2.55cm]
    \coordinate (A0) at (0,0); \coordinate (A1) at (1,0); \coordinate (A2) at (0.5,0.866);
    \draw[fb simplex] (A0) -- (A1) -- (A2) -- cycle;
    \foreach \c in {A0,A1,A2} \node[fb corner] at (\c) {};
    \node[fb lab,anchor=north east] at (A0) {$\alpha_0=1$};
    \node[fb lab,anchor=north west] at (A1) {$\alpha_1=1$};
    \node[fb lab,anchor=south] at (A2) {$\alpha_2=1$};
    \node[font=\small] at (0.5,0.3) {$\mathcal F>0$};
    \node[fb lab,anchor=north,align=center] at (0.5,-0.2) {Landau 解只在三个角上：\\零动量的无质量线};
    \node[font=\small,anchor=north,align=center] at (0.5,-0.75) {(b)\ 离壳（欧氏）：\\$\mathcal F=m^2\alpha_0(\alpha_1+\alpha_2)+Q^2\alpha_1\alpha_2$};
  \end{scope}
\end{tikzpicture}
\caption{单圈顶点的 Feynman 参数单形 $\{\alpha_0+\alpha_1+\alpha_2=1,\ \alpha_j\ge0\}$ 上的 Landau 集（示意；左为线与参数的记号）。
(a)~在壳无质量外线（$p_i^2=0$，$Q^2=2p_1\cdot p_2>0$）：$\mathcal F=-Q^2\alpha_1\alpha_2$，Landau 集为两条闭棱 $\alpha_2=0$（$k=xp_1$，共线于 $p_1$，红）与 $\alpha_1=0$（$k=-yp_2$，蓝），二者交于软角 $S$（$k=0$）；
另两个角是零动量夸克线（共线面的幂次压低端点）；棱 $\alpha_0=0$（虚线）只在 $Q^2=0$ 时是解；内部即硬区 $H$，没有钳夹（VF-REG-4）。
(b)~两条外线离壳（$p_i^2=-m^2$，欧氏非例外运动学）：$\mathcal F$ 只在单形的三个角上为零（零动量的无质量线，VF-REF-17）；欧氏区的红外有限性来自幂次计数，而不是因为没有 Landau 解（\cref{rem:ir-euclid-corners}）。}
\label{fig:ir-landau-simplex}
\end{figure}
%%% END fig:ir-landau-simplex


%%% BEGIN fig:ee-cutsum-disc
\begin{figure}[htbp]
\centering
% three-loop ladder graph of the current correlator (origin at its centre); #1 = cut path
\newcommand{\fbEEladder}[1]{%
  \draw[photon] (-1.85,0) -- (-1.25,0);
  \draw[photon] (1.25,0) -- (1.85,0);
  \draw[plain] (-1.25,0) to[out=62,in=180] (-0.62,0.72) -- (-0.42,0.72);
  \draw[plain] (-0.42,0.72) -- (0.42,0.72);
  \draw[quark=0.55] (0.42,0.72) -- (0.62,0.72) to[out=0,in=118] (1.25,0);
  \draw[plain] (1.25,0) to[out=-118,in=0] (0.62,-0.72) -- (0.42,-0.72);
  \draw[plain] (0.42,-0.72) -- (-0.42,-0.72);
  \draw[quark=0.55] (-0.42,-0.72) -- (-0.62,-0.72) to[out=180,in=-62] (-1.25,0);
  \draw[gluon] (-0.42,0.72) -- (-0.42,-0.72);
  \draw[gluon] (0.42,0.72) -- (0.42,-0.72);
  \foreach \p in {(-1.25,0),(1.25,0),(-0.42,0.72),(0.42,0.72),(-0.42,-0.72),(0.42,-0.72)} \node[vertex] at \p {};
  \draw[cut] #1;}
% reduced diagram of e+e- (origin at its centre); #1 = 1 draws the cut, 0 the forward box
\newcommand{\fbEEred}[1]{%
  \node[hard blob,circle,minimum size=0.42cm,inner sep=0pt] (H)  at (-1.13,0) {};
  \node[hard blob,circle,minimum size=0.42cm,inner sep=0pt] (Hd) at ( 1.13,0) {};
  \node[jet blob,ellipse,minimum width=1.42cm,minimum height=0.46cm,inner sep=0pt] (J1) at (0, 0.78) {};
  \node[jetB blob,ellipse,minimum width=1.42cm,minimum height=0.46cm,inner sep=0pt] (J2) at (0,-0.78) {};
  \node[soft blob,circle,minimum size=0.46cm,inner sep=0pt] (S) at (0,0) {};
  \draw[photon] (-1.5,0) -- (H.west);
  \draw[photon] (Hd.east) -- (1.5,0);
  \draw[quark] (H.70) -- (J1.188);
  \draw[quark] (J1.352) -- (Hd.110);
  \draw[quark] (Hd.250) -- (J2.8);
  \draw[quark] (J2.172) -- (H.290);
  \draw[gluon] (S.140) -- (-0.36,0.56);
  \draw[gluon] (S.40) -- (0.36,0.56);
  \draw[gluon] (S.220) -- (-0.36,-0.56);
  \draw[gluon] (S.320) -- (0.36,-0.56);
  \ifnum#1=1 \node[font=\scriptsize,inner sep=0pt,anchor=east] at ($(S.west)+(-0.04,0)$) {$S$};\else \node[font=\scriptsize] at (S) {$S$};\fi
  \node[font=\scriptsize] at (-0.47,0.78) {$J_1$};
  \node[font=\scriptsize] at (-0.47,-0.78) {$J_2$};
  \node[font=\scriptsize,anchor=north] at (H.south) {$H$};
  \node[font=\scriptsize,anchor=north] at (Hd.south) {$H^\dagger$};
  \ifnum#1=1 \draw[cut] (0,1.18) -- (0,-1.18);
  \else \draw[draw=gray!60!black,line width=0.6pt,dash pattern=on 2.5pt off 1.5pt,rounded corners=4pt]
          (-0.9,-1.14) rectangle (0.9,1.14);
  \fi}
\begin{tikzpicture}[fd,
  fb lab/.style={font=\footnotesize},
  fb note/.style={font=\scriptsize,text=red!55!black},
  fb bcut/.style={draw,line width=0.9pt,decorate,decoration={zigzag,segment length=4.2pt,amplitude=1.6pt}},
  fb contour/.style={draw=blue!60!black,line width=0.9pt},
  fb carrow/.style={postaction={decorate,decoration={markings,
      mark=at position #1 with {\arrow{Stealth[length=5.5pt,width=4.5pt]}}}}}]
  % ======================= (a) cut classes of one O(alpha_s^2) graph =======================
  \begin{scope}
    \fbEEladder{(0,1.05) -- (0,-1.05)}
    \node[fb lab] at (0,-1.38) {两粒子割线 $q\bar q$};
    \node[fb note] at (0,-1.76) {单类至多 $1/\varepsilon^4$ 极点};
  \end{scope}
  \begin{scope}[xshift=4.35cm]
    \fbEEladder{(-0.79,0.96) -- (0.09,-0.96)}
    \node[fb lab] at (0,-1.38) {三粒子割线 $q\bar qg$};
    \node[fb note] at (0,-1.76) {单类至多 $1/\varepsilon^4$ 极点};
  \end{scope}
  \begin{scope}[xshift=8.7cm]
    \fbEEladder{(-0.91,0.965) -- (0.91,-0.965)}
    \node[fb lab] at (0,-1.38) {四粒子割线 $q\bar qgg$};
    \node[fb note] at (0,-1.76) {单类至多 $1/\varepsilon^4$ 极点};
  \end{scope}
  % hand-made brace and the statement
  \draw[gray!45!black,line width=0.7pt] (-1.7,-2.0) to[out=-90,in=180] (-1.5,-2.2) -- (4.15,-2.2)
       to[out=0,in=90] (4.35,-2.4) to[out=90,in=180] (4.55,-2.2) -- (10.2,-2.2) to[out=0,in=-90] (10.4,-2.0);
  \node[font=\small] at (4.35,-2.75) {对全部割线求和（连同 $\MSbar$ 抵消项）：红外极点相消 $\;\Longrightarrow\;$ 有限};
  \node[font=\footnotesize,text=gray!30!black] at (4.35,-3.25) {（单态项 $\propto(\sum_qe_q)^2$ 自 $\order{\as^3}$ 起：Furry 定理）};
  \node[font=\small,anchor=west] at (-2.0,1.2) {(a)};
  % ======================= (b) cut sum = discontinuity of the forward function =======================
  \begin{scope}[yshift=-5.7cm]
    \node[font=\small,anchor=west] at (-2.0,1.45) {(b)};
    \node[font=\large] at (-1.55,0) {$\displaystyle\sum_{C}$};
    \begin{scope}[xshift=0.4cm]\fbEEred{1}\end{scope}
    \node[font=\normalsize] at (2.78,0) {$=\ \Disc_s$};
    \begin{scope}[xshift=5.25cm]\fbEEred{0}\end{scope}
    \node[fb lab,anchor=north] at (0.4,-1.28) {割线只穿过 $J_1$、$J_2$、$S$};
    \node[fb lab,anchor=north,align=center] at (5.25,-1.28) {前向（编时）组合 $F(s)$};
    % complex s plane
    \begin{scope}[xshift=10.0cm]
      \draw[lcaxis] (-1.95,0) -- (2.7,0) node[right,font=\footnotesize,inner sep=2pt] {$\Real s$};
      \draw[lcaxis] (0,-1.8) -- (0,1.95) node[left,font=\footnotesize] {$\Imag s$};
      \draw[fb bcut] (0,0) -- (2.45,0);
      \draw[fb contour,fb carrow=0.2,fb carrow=0.5,fb carrow=0.8] (5:1.45) arc[start angle=5,end angle=355,radius=1.45];
      \fill[blue!60!black] (5:1.45) circle[radius=1.5pt];
      \fill[blue!60!black] (-5:1.45) circle[radius=1.5pt];
      \node[font=\scriptsize,anchor=south west,inner sep=1pt] at ($(5:1.45)+(0.04,0.07)$) {$s+\ii0$};
      \node[font=\scriptsize,anchor=north west,inner sep=1pt] at ($(-5:1.45)+(0.04,-0.07)$) {$s-\ii0$};
      \node[font=\scriptsize,text=blue!55!black,anchor=south east,inner sep=1pt] at (132:1.5) {$|z|=s$：无钳夹};
      \node[font=\scriptsize,align=left,text=gray!25!black,anchor=north west,inner sep=1pt] at (1.62,-0.62) {钳夹（Landau\\奇点）只在\\实轴上};
      \node[font=\scriptsize,anchor=north] at (0.35,-1.95) {$F(s+\ii0)-F(s-\ii0)=-\displaystyle\int_{|z|=s}\!\dd z\,F'(z)$};
    \end{scope}
  \end{scope}
\end{tikzpicture}
\caption{割线求和与前向函数的不连续性（$e^+e^-$ 总截面；红色虚线为末态割线，左侧为振幅，右侧为共轭振幅）。
(a)~$\order{\as^2}$ 的一个三圈流关联函数图（夸克圈上两条胶子横档）的三类割线：两粒子 $q\bar q$、三粒子 $q\bar qg$ 与四粒子 $q\bar qgg$（每类只画一例，其余割线位置与镜像未画出）。
每类割图单独含有直到 $1/\varepsilon^4$ 的软、共线极点，只有对全部割线求和（连同 $\MSbar$ 抵消项）才有限（\cref{cor:ee-cuts}）；
两个夸克圈各连一个光子的单态项 $\propto(\sum_qe_q)^2$ 由 Furry 定理自 $\order{\as^3}$ 起才出现。
(b)~在钳夹面附近，割线只能穿过末态喷注 $J_1$、$J_2$ 与软子图 $S$，不能穿过 $H$、$H^\dagger$；对固定的 $H$、$H^\dagger$，由幺正性，对这些割线求和等于不加割线的前向（编时）组合 $F$（虚线框）的不连续性 $F(s+\ii0)-F(s-\ii0)$。
$F$ 是 $s$ 的解析函数，钳夹（Landau 奇点）只出现在实轴的割线上；由 Cauchy 定理，不连续性等于沿圆周 $|z|=s$ 的积分，圆周除两端外离实轴 $\order{s}$，那里没有钳夹，只剩欧氏区可计算的部分（围道公式，\cref{prop:ee-contour}）。}
\label{fig:ee-cutsum-disc}
\end{figure}
%%% END fig:ee-cutsum-disc


%%% BEGIN fig:dis-euclid-segment
\begin{figure}[htbp]
\centering
\begin{tikzpicture}[fd,
  fb bcut/.style={draw,line width=0.9pt,decorate,decoration={zigzag,segment length=4.2pt,amplitude=1.6pt}},
  fb seg/.style={draw=blue!65!black,line width=2.2pt,line cap=butt},
  fb lab/.style={font=\footnotesize},
  fb zoom/.style={draw=gray!55!black,line width=0.45pt}]
  \def\fbu{1.75}   % cm per unit of omega
  % ---------------- omega plane ----------------
  \fill[blue!5] (0,0) circle[radius=\fbu];
  \draw[gray!65!black,line width=0.7pt,dash pattern=on 1.2pt off 1.8pt] (0,0) circle[radius=\fbu];
  \draw[lcaxis] (-2.45*\fbu,0) -- (2.55*\fbu,0) node[above left=3pt and -2pt,font=\small] {$\Real\omega$};
  \draw[lcaxis] (0,-1.5*\fbu) -- (0,1.75*\fbu) node[left,font=\small] {$\Imag\omega$};
  \draw[fb bcut] (\fbu,0) -- (2.4*\fbu,0);
  \draw[fb bcut] (-\fbu,0) -- (-2.4*\fbu,0);
  \fill (\fbu,0) circle[radius=1.6pt]; \fill (-\fbu,0) circle[radius=1.6pt];
  \node[anchor=north west,font=\small,inner sep=1.5pt] at (\fbu+0.02,-0.08) {$1$};
  \node[anchor=north east,font=\small,inner sep=1.5pt] at (-\fbu-0.02,-0.08) {$-1$};
  \node[fb lab,anchor=north,align=center] at (1.9*\fbu,-0.16) {$s$ 道割线\\（物理区 $0<x\le1$）};
  \node[fb lab,anchor=north] at (-1.85*\fbu,-0.16) {$u$ 道割线};
  \node[fb lab,text=gray!35!black,anchor=north east] at (-0.72*\fbu,-0.72*\fbu) {$|\omega|=1$};
  \node[fb lab,align=left,text=blue!45!black,anchor=south east] (AN) at (-1.05,1.82) {单位圆盘 $|\omega|<1$ 内\\$T_2$ 解析（\cref{prop:dis-dispersion-np}）};
  \draw[blue!45!black,line width=0.4pt] ($(AN.south east)+(-0.25,0.02)$) -- (-0.72,0.9);
  % the Euclidean segment I_Q (drawn to scale for sqrt(3)M/Q = 0.08)
  \draw[fb seg] (0,-0.08*\fbu) -- (0,0.08*\fbu);
  % magnifier: small circle at the origin, external tangents to the inset
  \coordinate (Z) at (5.7,3.45);
  \draw[fb zoom] (0,0) circle[radius=0.49];
  \draw[fb zoom] (-0.320,0.371) -- (4.655,4.662);
  \draw[fb zoom] (0.180,-0.456) -- (6.289,1.962);
  \begin{scope}[shift={(Z)}]
    \begin{scope}
      \clip (0,0) circle[radius=1.6];
      \fill[blue!5] (-1.7,-1.7) rectangle (1.7,1.7);
      % the origin of omega sits at (-0.5,0) inside the inset
      \draw[gray!60!black,line width=0.5pt] (-1.7,0) -- (1.7,0);
      \draw[gray!60!black,line width=0.5pt] (-0.5,-1.7) -- (-0.5,1.7);
    \end{scope}
    \draw[fb zoom,line width=0.7pt] (0,0) circle[radius=1.6];
    \draw[fb seg] (-0.5,-0.75) -- (-0.5,0.75);
    \draw[blue!65!black,line width=0.6pt] (-0.61,0.75) -- (-0.39,0.75) (-0.61,-0.75) -- (-0.39,-0.75);
    \node[font=\scriptsize,anchor=west,inner sep=2pt] at (-0.38,0.75) {$\ii\sqrt3\,M/Q$};
    \node[font=\scriptsize,anchor=west,inner sep=2pt] at (-0.38,-0.75) {$-\ii\sqrt3\,M/Q$};
    \node[font=\small,text=blue!55!black,anchor=west] at (-0.36,0.3) {$\mathcal I_Q$};
    \draw[red!60!black,line width=0.7pt,-{Stealth[length=4pt,width=3.4pt]}] (-0.9,1.27) -- (-0.9,0.82);
    \draw[red!60!black,line width=0.7pt,-{Stealth[length=4pt,width=3.4pt]}] (-0.9,-1.27) -- (-0.9,-0.82);
  \end{scope}
  \node[font=\footnotesize,anchor=west,align=left] at ($(Z)+(1.72,0)$) {{\color{gray!25!black}原点附近放大}\\[3pt]{\color{red!55!black}$Q\to\infty$：}\\{\color{red!55!black}线段缩向原点}};
\end{tikzpicture}
\caption{前向 Compton 振幅 $T_2$ 在复 $\omega=1/x$ 平面上的欧氏点。靶静止系中，欧氏动量 $q=(\ii q^4,\mathbf q)$（$|\mathbf q|\ge Q/2$）给出 $\omega=2\ii Mq^4/Q^2$，
恰好取遍虚轴上的短线段 $\mathcal I_Q=\{\ii y:\,|y|\le\sqrt3M/Q\}$（蓝色粗线；右上为原点附近的放大），$Q\to\infty$ 时它缩向原点（VF-REF-1）。
$T_2$ 在单位圆盘 $|\omega|<1$ 内解析，割线 $(-\infty,-1]$ 与 $[1,\infty)$ 分别为 $u$ 道与 $s$ 道（物理区）的割线。
欧氏 OPE 只在长度为 $\order{M/Q}$ 的线段 $\mathcal I_Q$ 上提供信息；由此得到固定 $n$ 的矩还需要 $F_2\ge0$ 与有限减除的色散关系（\cref{thm:dis-fixed-n}），
而 $n<n_0$ 的低阶矩还需要增长条件（\cref{prop:dis-low-moments}）。}
\label{fig:dis-euclid-segment}
\end{figure}
%%% END fig:dis-euclid-segment


%%% BEGIN fig:method-flowchart
\begin{figure}[htbp]
\centering
\newcommand{\fbref}[1]{\mbox{\cref{#1}}}
\begin{tikzpicture}[fd,font=\footnotesize,
  fb num/.style={circle,fill=gray!65!black,text=white,font=\scriptsize\bfseries,inner sep=0pt,minimum size=0.42cm},
  fb title/.style={font=\small\bfseries,anchor=west,inner sep=1pt},
  fb stitle/.style={font=\footnotesize\bfseries,anchor=west,inner sep=1pt},
  fb cell/.style={rectangle,rounded corners=1.5pt,line width=0.5pt,text width=3.38cm,minimum height=1.0cm,
                  align=left,inner sep=3pt,font=\footnotesize},
  fb R/.style={fb cell,fill=OliveGreen!10,draw=OliveGreen!65!black},
  fb Rc/.style={fb R,dash pattern=on 2.4pt off 1.4pt},
  fb A/.style={fb cell,fill=RoyalBlue!8,draw=RoyalBlue!65!black},
  fb frame/.style={draw=gray!55!black,line width=0.65pt,rounded corners=3pt,fill=white,inner sep=3pt},
  fb sub/.style={draw=gray!45,line width=0.45pt,rounded corners=2pt,fill=gray!4,inner sep=2pt},
  fb htext/.style={text width=6.45cm,align=left,inner sep=1pt,font=\footnotesize},
  fb H/.style={rectangle,rounded corners=2pt,draw=orange!70!black,line width=0.55pt,dash pattern=on 2.4pt off 1.4pt,
               fill=orange!5,inner sep=3pt},
  fb link/.style={draw=orange!70!black,line width=0.5pt,dash pattern=on 2.4pt off 1.4pt},
  fb flow/.style={draw=gray!55!black,line width=1.1pt,-{Stealth[length=6pt,width=5pt]}}]
  % one step: #1 name, #2 y of its title, #3 number, #4 title, #5 R cell, #6 A cell, #7 R style
  \newcommand{\fbstep}[7]{%
    \node[fb num] (#1n) at (0.28,#2) {#3};
    \node[fb title] (#1t) at ($(#1n.east)+(0.06,0)$) {#4};
    \node[#7,anchor=north west] (#1r) at (0.12,#2-0.3) {#5};
    \node[fb A,anchor=north west] (#1a) at ($(#1r.north east)+(0.12,0)$) {#6};
    \begin{scope}[on background layer]\node[fb frame,fit=(#1n)(#1t)(#1r)(#1a)] (#1) {};\end{scope}}
  % ---------------- steps 1--4 ----------------
  \fbstep{s1}{0}{1}{区域：Landau 方程与幂次计数}
    {幂次计数（标度陈述）\\\fbref{prop:ir-power-counting}}
    {领头区域与约化图\\\fbref{thm:ir-leading-regions}}{fb R}
  \fbstep{s2}{-2.1}{2}{近似：共线与软近似子（沿用原有 $\ii0$）}
    {$\ii0$ 相容性\\\fbref{lem:method-i0-compat}}
    {近似精度\\\fbref{lem:method-collinear-accuracy,lem:method-soft-accuracy}}{fb R}
  \fbstep{s3}{-4.2}{3}{Grammer--Yennie 分解\enspace$g_{\mu\nu}=K_{\mu\nu}+G_{\mu\nu}$}
    {分解恒等式\\\fbref{lem:method-gy}}
    {G 部分幂次压低\\\fbref{prop:method-gy-suppression}}{fb R}
  \fbstep{s4}{-6.3}{4}{Ward 恒等式 $\to$ Wilson 线}
    {单个与树图两个 K 胶子\\\fbref{prop:method-ward-one,prop:method-ward-tree}}
    {多胶子（含鬼场）\\\fbref{lem:method-ward-multi}}{fb R}
  % ---------------- step 5: Glauber region, unitarity, subtraction ----------------
  \node[fb num] (s5n) at (0.28,-8.4) {5};
  \node[fb title] (s5t) at ($(s5n.east)+(0.06,0)$) {Glauber 区、幺正性与减除};
  \foreach \k/\y/\ttl/\rc/\ac/\rs in {%
      a/-8.85/{5a\enspace Glauber 区：沿 $e_G=n_B-n_A$ 同时形变}/{单胶子交换的极点结构\\\fbref{prop:method-glauber-oneloop}}/{围道形变判据\\\fbref{prop:method-glauber-deformation}}/{fb R},
      b/-10.75/{5b\enspace 幺正性：末态相互作用与软因子}/{软因子的幺正性\\\fbref{lem:method-soft-unitarity}}/{末态相互作用相消\\\fbref{prop:method-fsi}}/{fb R},
      c/-12.65/{5c\enspace 减除（Collins）：$\Gamma=\sum_RC_R\Gamma+r$}/{链展开与乘积公式\\\fbref{prop:method-chains,prop:method-product}}/{减除定理\\\fbref{thm:method-subtraction}}/{fb R}} {
    \node[fb stitle] (g5\k t) at (0.22,\y) {\ttl};
    \node[\rs,anchor=north west,text width=3.3cm] (g5\k r) at (0.22,\y-0.24) {\rc};
    \node[fb A,anchor=north west,text width=3.3cm] (g5\k a) at ($(g5\k r.north east)+(0.12,0)$) {\ac};
    \begin{scope}[on background layer]\node[fb sub,fit=(g5\k t)(g5\k r)(g5\k a)] (s5\k) {};\end{scope}
  }
  \begin{scope}[on background layer]
    \node[fb frame,fit=(s5n)(s5t)(s5a)(s5b)(s5c)] (s5) {};
  \end{scope}
  % ---------------- step 6 ----------------
  \node[fb num] (s6n) at (0.28,-15.05) {6};
  \node[fb title] (s6t) at ($(s6n.east)+(0.06,0)$) {因子化形式};
  \node[fb A,anchor=north west,text width=7.09cm,minimum height=0.62cm,align=center] (s6a) at (0.12,-15.35)
    {$\sigma=\sum_{ij}H_{ij}\otimes f_{i/A}\otimes f_{j/B}\,[\otimes S]+\order{\lambda^p}$};
  \begin{scope}[on background layer]\node[fb frame,fit=(s6n)(s6t)(s6a)] (s6) {};\end{scope}
  % ---------------- flow arrows ----------------
  \foreach \u/\v in {s1/s2,s2/s3,s3/s4,s4/s5,s5/s6} \draw[fb flow] (\u.south) -- (\v.north);
  % ---------------- hypotheses (H1)--(H5) with the numbered assumptions ----------------
  \node[fb htext,anchor=west] (h1t) at (8.2,0|-s1) {\textbf{(H1)} 领头区域；积分范数下的幂次计数\\\fbref{asm:appN-envelope,asm:appN-separation}};
  \node[fb htext,anchor=west] (h2t) at (8.2,0|-s3) {\textbf{(H2)} $H$--$J$ 界面上共线近似成立，K 胶子退耦为 Wilson 线\\\fbref{asm:appN-ward-product}\\[1pt]
      {\color{gray!35!black}（步骤 2--4 中的软部分属于 (H3)）}};
  \node[fb htext,anchor=north west] (h3t) at ($(8.2,0|-s5a)+(0,0.66)$) {\textbf{(H3)} Glauber 区之后：软近似与软胶子退耦\\\fbref{asm:appO-mixing,asm:appO-deformation,asm:appO-rewritten-cs}};
  \node[fb Rc,anchor=north west,text width=6.3cm,minimum height=0.48cm] (h3c) at ($(h3t.south west)+(0,-0.07)$) {DY 的精确形式：\fbref{thm:appO-main}（$R_{\rm cond}$）};
  \node[fb htext,anchor=west] (h4t) at (8.2,0|-s5b) {\textbf{(H4)} 末态相互作用由幺正性相消；软因子为 $1$ 或 Wilson 线真空矩阵元\\\fbref{asm:appN-cutsum,asm:appO-lcopt}};
  \node[fb htext,anchor=north west] (h5t) at ($(8.2,0|-s5c)+(0,0.66)$) {\textbf{(H5)} 减除项 = 重整化的算符矩阵元\\\fbref{asm:appN-collins,asm:appN-values,asm:appN-uv,asm:appN-multren}};
  \node[fb Rc,anchor=north west,text width=6.3cm,minimum height=0.48cm] (h5c) at ($(h5t.south west)+(0,-0.07)$) {精确演绎：\fbref{thm:appN-subtraction}（$R_{\rm cond}$）};
  \begin{scope}[on background layer]
    \node[fb H,fit=(h1t)] (h1) {};
    \node[fb H,fit=(h2t)] (h2) {};
    \node[fb H,fit=(h3t)(h3c)] (h3) {};
    \node[fb H,fit=(h4t)] (h4) {};
    \node[fb H,fit=(h5t)(h5c)] (h5) {};
  \end{scope}
  \draw[fb link] (s1.east) -- (h1.west|-s1.east);
  \draw[fb link] (s2.east) -- ++(0.25,0) |- (h2.west);
  \draw[fb link] (s3.east) -- (h2.west|-s3.east);
  \draw[fb link] (s4.east) -- ++(0.25,0) |- (h2.west);
  \draw[fb link] (s5a.east) -- (h3.west|-s5a.east);
  \draw[fb link] (s5b.east) -- (h4.west|-s5b.east);
  \draw[fb link] (s5c.east) -- (h5.west|-s5c.east);
  \node[font=\footnotesize,anchor=west,align=left] at (8.2,0|-s6) {$\Longleftarrow$\ \mbox{(H1)--(H5)} 成立时的结论：\\\fbref{thm:method-template}（模板定理），全阶论证 [A]};
  % ---------------- legend ----------------
  \node[fb R,text width=0.3cm,minimum height=0.3cm,inner sep=0pt,anchor=west] (l1) at (0.1,-16.55) {};
  \node[anchor=west] (l1t) at (l1.east) {[R]（在所用层次上严格；含数学引理）};
  \node[fb A,text width=0.3cm,minimum height=0.3cm,inner sep=0pt,anchor=west] (l2) at ($(l1t.east)+(0.35,0)$) {};
  \node[anchor=west] (l2t) at (l2.east) {[A] 全阶论证};
  \node[fb Rc,text width=0.3cm,minimum height=0.3cm,inner sep=0pt,anchor=west] (l3) at ($(l2t.east)+(0.35,0)$) {};
  \node[anchor=west] (l3t) at (l3.east) {条件定理 $R_{\rm cond}$};
  \node[fb H,minimum width=0.36cm,minimum height=0.32cm,inner sep=0pt,anchor=west] (l4) at ($(l3t.east)+(0.35,0)$) {};
  \node[anchor=west] (l4t) at (l4.east) {编号假设};
\end{tikzpicture}
\caption{全阶因子化证明的六个步骤（左），它们与\cref{thm:method-template}（模板定理）的假设 \mbox{(H1)--(H5)} 的对应，以及精确形式中的编号假设（右；\cref{app:appN}与\cref{app:appO}，对应关系即\cref{tab:method-hypotheses}）。
每个步骤中，绿色框为在所用层次上已严格证明的陈述（[R] 或数学引理），蓝色框为全阶论证 [A]，绿色虚线框为条件定理 $R_{\rm cond}$，即\cref{app:appN}中的减除演绎（\cref{thm:appN-subtraction}）与\cref{app:appO}中单举色单态 Drell--Yan 过程的 Glauber 相消（\cref{thm:appO-main}）。
步骤 4 中，$A$ 共线的 K 胶子化为沿 $n_B$ 的 Wilson 线（$B$ 共线的则沿 $n_A$），软胶子化为沿各喷注方向的 Wilson 线；离开 Glauber 区须沿 $e_G=n_B-n_A$ 同时形变两个光锥分量（\cref{fig:method-glauber-direction}），软因子的幺正性只在处理完 Glauber 区之后使用。}
\label{fig:method-flowchart}
\end{figure}
%%% END fig:method-flowchart


%%% BEGIN fig:method-hasse
\begin{figure}[htbp]
\centering
\begin{tikzpicture}[fd,
  fb J1/.style={draw=blue!60!black},
  fb J2/.style={draw=orange!80!black},
  fb S/.style={draw=green!45!black},
  fb Hl/.style={draw=black,line width=1.6pt},
  fb m/.style={font=\scriptsize,inner sep=2.6pt,sloped},
  fb node/.style={draw=gray!60!black,fill=white,rounded corners=2.5pt,line width=0.6pt,inner sep=2.5pt,
                  minimum width=1.5cm,minimum height=0.55cm,font=\small},
  fb edge/.style={draw=gray!50!black,line width=0.7pt},
  fb hl/.style={draw=red!60!black,line width=1.3pt}]
  \newcommand{\fbH}{H}
  \newcommand{\fbA}{{\color{blue!60!black}C_1}}
  \newcommand{\fbB}{{\color{orange!80!black}C_2}}
  \newcommand{\fbS}{{\color{green!45!black}S}}
  % ======================= left: two-loop planar vertex, coloured for the region (C1,S) =======================
  \begin{scope}
    \coordinate (V) at (0,0);
    \coordinate (Q1) at (28:1.9); \coordinate (Q2) at (28:3.4); \coordinate (E1) at (28:4.55);
    \coordinate (A1) at (-28:1.9); \coordinate (A2) at (-28:3.4); \coordinate (E2) at (-28:4.55);
    \draw[photon] (-1.3,0) -- (V);
    \draw[quark=0.82,fb Hl] (A1) -- node[fb m,below,pos=0.4,text=black] {$p_2+k_1+k_2$} (V);
    \draw[quark=0.2,fb J1] (V) -- node[fb m,above,pos=0.6,text=blue!50!black] {$p_1-k_1-k_2$} (Q1);
    \draw[quark=0.22,fb J1] (Q1) -- node[fb m,above,pos=0.62,text=blue!50!black] {$p_1-k_2$} (Q2);
    \draw[quark] (Q2) -- (E1);
    \draw[quark=0.8,fb J2] (A2) -- node[fb m,below,pos=0.38,text=orange!60!black] {$p_2+k_2$} (A1);
    \draw[quark] (E2) -- (A2);
    \draw[gluon,fb J1] (A1) -- (Q1);
    \draw[gluon,fb S] (A2) -- (Q2);
    \foreach \p in {V,Q1,Q2,A1,A2} \node[vertex] at (\p) {};
    \path[mom={$q$},mom sep=0.28cm] (-1.3,0) -- (V);
    \path[mom={$p_1$},mom pos=0.6,mom length=0.45cm] (Q2) -- (E1);
    \path[mom={$p_2$},mom pos=0.6,mom length=0.45cm] (A2) -- (E2);
    \path[mom'={\scriptsize$k_1$},mom length=0.42cm,mom sep=0.3cm] (A1) -- (Q1);
    \path[mom'={\scriptsize$k_2$},mom length=0.42cm,mom sep=0.3cm] (A2) -- (Q2);
    \node[font=\footnotesize,align=center] at (1.75,-2.75) {例：区域 $(R(k_1),R(k_2))=(\fbA,\fbS)$\\[1pt]
      {\color{blue!60!black}$J_1$}\quad{\color{orange!80!black}$J_2$}\quad{\color{green!45!black}$S$}\quad 粗线：$H$};
  \end{scope}
  % ======================= right: Hasse diagram of the nine leading regions =======================
  \begin{scope}[xshift=8.6cm,yshift=-2.35cm]
    \node[fb node] (SS) at (0,0) {$(\fbS,\fbS)$};
    \node[fb node] (AS) at (-1.45,1.15) {$(\fbA,\fbS)$};
    \node[fb node] (BS) at (1.45,1.15) {$(\fbB,\fbS)$};
    \node[fb node] (AA) at (-2.75,2.3) {$(\fbA,\fbA)$};
    \node[fb node] (HS) at (0,2.3) {$(\fbH,\fbS)$};
    \node[fb node] (BB) at (2.75,2.3) {$(\fbB,\fbB)$};
    \node[fb node] (HA) at (-1.45,3.45) {$(\fbH,\fbA)$};
    \node[fb node] (HB) at (1.45,3.45) {$(\fbH,\fbB)$};
    \node[fb node] (HH) at (0,4.6) {$(\fbH,\fbH)$};
    \foreach \u/\v in {SS/AS,SS/BS,AS/AA,AS/HS,BS/BB,BS/HS,HS/HA,HS/HB,BB/HB,HA/HH,HB/HH}
      \draw[fb edge] (\u) -- (\v);
    \draw[fb hl] (AA) -- (HA);
    \node[font=\scriptsize,text=red!55!black,align=right,anchor=east] at (-2.35,3.5) {逐线比较：\\$C_1\le H$，$C_1\le C_1$};
    % beta_R column
    \node[font=\footnotesize,text=gray!30!black] at (4.25,5.35) {$\beta_R$};
    \foreach \y/\b in {0/{1/81},1.15/{1/27},2.3/{1/9},3.45/{1/3}}
      \node[font=\footnotesize,text=gray!30!black] at (4.25,\y) {$\b$};
    \node[font=\scriptsize,text=gray!30!black] at (4.25,4.6) {$N=\Omega$};
    % pinch sets grow upwards
    \draw[lcaxis,gray!50!black,line width=0.7pt] (5.05,-0.2) -- (5.05,4.8);
    \node[font=\footnotesize,text=gray!30!black,rotate=90,anchor=north] at (5.12,2.3) {钳夹集增大};
    \node[font=\footnotesize,anchor=south] at (0,5.1) {$(R(k_1),R(k_2))$：内胶子与外胶子的模式};
  \end{scope}
\end{tikzpicture}
\caption{两圈平面顶点九个领头区域的 Hasse 图。左：内胶子 $k_1$、外胶子 $k_2$ 的动量走向（胶子动量由反夸克线流向夸克线），线的颜色给出区域 $(C_1,S)$ 的逐线模式：
$J_1=\{k_1,\,p_1-k_2,\,p_1-k_1-k_2\}$、$J_2=\{p_2+k_2\}$、$S=\{k_2\}$，硬线 $p_2+k_1+k_2$ 与光子顶点组成 $H$。
右：16 种模式组合中恰有 9 个领头钳夹，按逐线偏序（$S<C_i<H$，$C_1\wedge C_2=S$）排成 Hasse 图（12 条覆盖关系，自下而上钳夹集增大），其中 $(C_1,C_1)<(H,C_1)$（红线）不易一眼看出；
这一集合对逐线交运算封闭，例如 $(C_1,C_1)\wedge(H,S)=(C_1,S)$、$(H,C_1)\wedge(H,C_2)=(H,S)$、$(C_1,C_1)\wedge(C_2,C_2)=(S,S)$。
右列为各层的邻域指数 $\beta_R=3^{h(R)-h_*}$（$h_*=4$；最大区域的邻域为整个积分域 $\Omega$，\cref{def:method-neighbourhoods}）。}
\label{fig:method-hasse}
\end{figure}
%%% END fig:method-hasse


%%% BEGIN fig:method-nesting
\begin{figure}[htbp]
\centering
\begin{tikzpicture}[fd,
  fb lab/.style={font=\footnotesize},
  fb ellA/.style={rotate around={22:(-0.95,0)}},
  fb ellB/.style={rotate around={-22:(0.95,0)}},
  fb step/.style={draw=gray!40!black,line width=0.75pt,-{Stealth[length=5pt,width=4.2pt]}},
  fb stepl/.style={font=\footnotesize,inner sep=1.5pt,fill=white,rounded corners=1.5pt}]
  % ---------------- Venn-like diagram of the pinch sets (schematic) ----------------
  \draw[fill=gray!10,draw=gray!55!black,line width=0.7pt,rounded corners=10pt] (-3.6,-2.35) rectangle (3.6,2.35);
  \fill[blue!12,fb ellA] (-0.95,0) ellipse[x radius=2.05,y radius=1.08];
  \fill[orange!17,fb ellB] (0.95,0) ellipse[x radius=2.05,y radius=1.08];
  % pi_S = pi_A cap pi_B: the lens common to both ellipses
  \begin{scope}
    \clip[fb ellA] (-0.95,0) ellipse[x radius=2.05,y radius=1.08];
    \fill[green!22,fb ellB] (0.95,0) ellipse[x radius=2.05,y radius=1.08];
  \end{scope}
  \draw[blue!55!black,line width=0.7pt,fb ellA] (-0.95,0) ellipse[x radius=2.05,y radius=1.08];
  \draw[orange!75!black,line width=0.7pt,fb ellB] (0.95,0) ellipse[x radius=2.05,y radius=1.08];
  \node[font=\small,text=green!30!black] at (0,0.22) {$\pi_S$};
  \node[font=\small,text=blue!50!black] at (-2.05,-0.25) {$\pi_A$};
  \node[font=\small,text=orange!60!black] at (2.05,-0.25) {$\pi_B$};
  \node[font=\small,anchor=north west] at (-3.5,2.25) {$\pi_H=\Omega$};
  \node[fb lab,text=gray!30!black,anchor=south] at (0,-2.3) {$\pi_S=\pi_A\cap\pi_B$（\cref{lem:method-meets}）};
  % order of the approximations: smallest region first
  \draw[fb step] (0.1,1.95) node[fb stepl,anchor=south] {\textcircled{\scriptsize 1}\,$T_S$} to[out=-90,in=90] (0.04,0.62);
  \draw[fb step] (-2.9,-1.65) node[fb stepl,anchor=north] {\textcircled{\scriptsize 2}\,$T_A$} to[out=60,in=-150] (-1.85,-0.75);
  \draw[fb step] (2.9,-1.65) node[fb stepl,anchor=north] {\textcircled{\scriptsize 2}\,$T_B$} to[out=120,in=-30] (1.85,-0.75);
  \draw[fb step] (2.75,1.85) node[fb stepl,anchor=south west,xshift=-0.15cm] {\textcircled{\scriptsize 3}\,$T_H$} to[out=-120,in=60] (2.4,1.2);
  % ---------------- subtraction terms ----------------
  \begin{scope}[xshift=4.15cm]
    \node[anchor=north west,align=left,font=\small] (F) at (0,2.2) {%
      $C_S\Gamma=T_S\,\Gamma$\\[4pt]
      $C_A\Gamma=T_A(1-T_S)\,\Gamma$，\quad$C_B\Gamma=T_B(1-T_S)\,\Gamma$\\[4pt]
      $C_H\Gamma=T_H\big(\Gamma-C_S\Gamma-C_A\Gamma-C_B\Gamma\big)$\\[4pt]
      $\Gamma=\sum_{R}C_R\Gamma+r$};
    \node[anchor=north west,align=left,font=\footnotesize,text width=6.6cm,draw=red!50!black,
          line width=0.5pt,rounded corners=2pt,fill=red!3,inner sep=4pt] at ($(F.south west)+(0.05,-0.25)$)
      {只有分组表达式被压低（\cref{ex:method-naive-induction}）：单个链项一般不可积（例如离壳三角图中的 $T_{C_1}\Gamma$ 对任何 $\varepsilon$ 都不绝对可积），被压低并满足 \mbox{(CS1)--(CS3)} 的是分组表达式 $C_R\Gamma$ 与余项 $r$。};
  \end{scope}
\end{tikzpicture}
\caption{区域的偏序与减除（示意；钳夹集画成区域）。单圈顶点的钳夹集彼此嵌套：软钳夹集 $\pi_S$（绿色）同时含于两个共线钳夹集 $\pi_A$（蓝）与 $\pi_B$（橙），并且恰为二者之交 $\pi_S=\pi_A\cap\pi_B$，即逐线交 $A\wedge B=S$；所有钳夹集都含于硬区 $\pi_H=\Omega$。
减除自最小区域（软）开始，依次为 \textcircled{\scriptsize 1}~$T_S$、\textcircled{\scriptsize 2}~$T_A,T_B$、\textcircled{\scriptsize 3}~$T_H$：较大区域的近似子作用于减去较小区域贡献之后的被积函数，$C_R\Gamma\equiv T_R\big(\Gamma-\sum_{R'<R}C_{R'}\Gamma\big)$（\cref{def:method-subtraction}）。
单个链项一般不可积，只有分组表达式满足 \mbox{(CS1)--(CS3)}。}
\label{fig:method-nesting}
\end{figure}
%%% END fig:method-nesting


%%% BEGIN fig:method-neighbourhoods
\begin{figure}[htbp]
\centering
\begin{tikzpicture}[fd,
  fb lab/.style={font=\footnotesize},
  fb dim/.style={draw=black!75,line width=0.5pt,{Stealth[length=3.6pt,width=3pt]}-{Stealth[length=3.6pt,width=3pt]}},
  fb lead/.style={draw=black!60,line width=0.4pt}]
  \def\fbw{0.42}   % half-width lambda^{beta_C} of the collinear tubes
  \def\fbr{1.15}   % radius lambda^{beta_S} of the soft ball (beta_S < beta_C: fatter)
  % ---------------- the domain Omega and the dominance sets ----------------
  \begin{scope}
    \clip[rounded corners=9pt] (-3.75,-2.45) rectangle (3.75,2.45);
    \fill[gray!12] (-3.8,-2.5) rectangle (3.8,2.5);
    \fill[blue!15] (-3.8,-\fbw) rectangle (3.8,\fbw);
    \fill[orange!20] (-\fbw,-2.5) rectangle (\fbw,2.5);
    \draw[blue!55!black,line width=0.5pt,dash pattern=on 2pt off 1.5pt] (-3.8,-\fbw) -- (3.8,-\fbw) (-3.8,\fbw) -- (3.8,\fbw);
    \draw[orange!75!black,line width=0.5pt,dash pattern=on 2pt off 1.5pt] (-\fbw,-2.5) -- (-\fbw,2.5) (\fbw,-2.5) -- (\fbw,2.5);
    \fill[green!20] (0,0) circle[radius=\fbr];
    \draw[green!45!black,line width=0.6pt] (0,0) circle[radius=\fbr];
    \draw[blue!60!black,line width=1.1pt] (-3.8,0) -- (3.8,0);
    \draw[orange!80!black,line width=1.1pt] (0,-2.5) -- (0,2.5);
  \end{scope}
  \draw[gray!55!black,line width=0.7pt,rounded corners=9pt] (-3.75,-2.45) rectangle (3.75,2.45);
  \fill (0,0) circle[radius=1.8pt];
  % axes of the section k_T=0
  \draw[lcaxis,blue!60!black] (3.75,0) -- (4.1,0) node[right=0pt,font=\footnotesize,text=black] {$k^+$};
  \draw[lcaxis,orange!80!black] (0,2.45) -- (0,2.8) node[above=0pt,font=\footnotesize,text=black] {$k^-$};
  % pinch sets
  \node[fb lab,text=blue!50!black,anchor=south west] at (-3.62,\fbw+0.02) {$\pi_{C_1}$（$k\parallel p_1$）};
  \draw[fb lead,blue!50!black] (-3.15,\fbw+0.07) -- (-3.0,0.04);
  \node[fb lab,text=orange!60!black,anchor=north west] at (\fbw+0.06,2.36) {$\pi_{C_2}$（$k\parallel p_2$）};
  \draw[fb lead,orange!60!black] (\fbw+0.1,2.12) -- (0.04,1.95);
  \node[fb lab,anchor=south west,inner sep=1.5pt] at (0.05,0.03) {$\pi_S$};
  % dominance sets
  \node[font=\small,text=green!30!black] at (-0.66,-0.7) {$U_S$};
  \node[font=\small,text=blue!45!black] at (2.35,-0.21) {$U_{C_1}$};
  \node[font=\small,text=blue!45!black] at (-2.35,-0.21) {$U_{C_1}$};
  \node[font=\small,text=orange!55!black,fill=orange!20,inner sep=1pt] at (0,-1.78) {$U_{C_2}$};
  \node[font=\small,text=orange!55!black,fill=orange!20,inner sep=1pt] at (0,1.55) {$U_{C_2}$};
  \node[font=\small,text=gray!35!black] at (-2.35,1.55) {$U_H$};
  \node[font=\small,text=gray!35!black] at (2.35,-1.55) {$U_H$};
  \node[fb lab,anchor=south east,text=gray!30!black] at (3.68,-2.4) {$\Omega$};
  % dimensions
  \draw[fb dim] (0,0) -- (-45:\fbr);
  \draw[fb lead] (-45:0.62) -- (1.45,-1.12);
  \node[fb lab,anchor=west,inner sep=1pt] at (1.45,-1.12) {$\lambda^{\beta_S}$};
  \draw[fb dim] (3.05,0) -- (3.05,\fbw);
  \draw[fb lead] (3.05,0.3) -- (3.3,0.75);
  \node[fb lab,anchor=south,inner sep=1pt] at (3.3,0.75) {$\lambda^{\beta_C}$};
  % ---------------- definitions ----------------
  \node[anchor=north west,align=left,font=\footnotesize,text width=6.3cm] at (4.75,2.42) {%
    $N_R=\{\sigma_{R(j)}(l_j)\le\lambda^{\beta_R}\ \ \forall j:\ R(j)\neq H\}$\\[2pt]
    $\sigma_A=|l^-|+|\mathbf l_T|$，\ $\sigma_B=|l^+|+|\mathbf l_T|$，\ $\sigma_S=\|l\|$\\[6pt]
    $U_R=N_R\setminus\bigcup_{R'<R}N_{R'}$：\\[2pt]
    \quad$U_S=N_S$，\ \ $U_{C_i}=N_{C_i}\setminus N_S$，\\[2pt]
    \quad$U_H=\Omega\setminus(N_S\cup N_{C_1}\cup N_{C_2})$\\[6pt]
    $\beta_R=3^{h(R)-h_*}$：较小区域的邻域较“胖”\\[2pt]
    （本例 $\beta_S=\frac19<\beta_C=\frac13$，$\lambda^{\beta_S}>\lambda^{\beta_C}$）\\[6pt]
    $\bigcup_RU_R=\Omega$（\cref{lem:method-covering}）};
\end{tikzpicture}
\caption{依赖于 $\lambda$ 的邻域 $N_R=\{\sigma_{R(j)}(l_j)\le\lambda^{\beta_R}\}$ 与支配集 $U_R=N_R\setminus\bigcup_{R'<R}N_{R'}$（\cref{def:method-neighbourhoods}；示意：单圈离壳三角图的胶子动量 $k$ 在 $\mathbf k_T=0$ 处的 $(k^+,k^-)$ 截面，以 $Q$ 为单位）。
共线钳夹集 $\pi_{C_1}$、$\pi_{C_2}$ 相交于软钳夹点 $\pi_S$（$k=0$）；$N_{C_i}$ 是半宽 $\lambda^{\beta_C}$ 的管道，$N_S$ 是半径 $\lambda^{\beta_S}$ 的球，
由于 $\beta_S<\beta_C$ 而比管道“胖”。绿、蓝、橙、灰色区域分别为 $U_S$、$U_{C_1}$、$U_{C_2}$、$U_H$，它们覆盖整个积分域（\cref{lem:method-covering}）。}
\label{fig:method-neighbourhoods}
\end{figure}
%%% END fig:method-neighbourhoods


%%% BEGIN fig:method-kg
\begin{figure}[htbp]
\centering
% collinear blob A, hard blob, photon and outgoing struck quark (origin at the centre of A); #1 = name of the hard blob
\newcommand{\fbKGframe}[1]{%
  \node[jet blob,ellipse,minimum width=1.75cm,minimum height=0.66cm,inner sep=0pt] (A) at (0,0) {};
  \node[font=\small] at (-0.42,0) {$A$};
  \node[hard blob,ellipse,minimum width=1.6cm,minimum height=0.66cm,inner sep=0pt] (H) at (0,2.25) {};
  \node[font=\small] at (H) {$#1$};
  \draw[photon] (-1.25,3.05) -- (H.150);
  \draw[quark] (H.20) -- (1.25,3.0);}
\begin{tikzpicture}[fd,
  fb lab/.style={font=\footnotesize},
  fb idx/.style={font=\scriptsize,inner sep=1pt},
  fb op/.style={font=\large},
  fb box/.style={draw=gray!60!black,line width=0.5pt,rounded corners=3pt,fill=gray!4,inner sep=4pt}]
  % ======================= row 1: g = K + G =======================
  % full propagator numerator
  \begin{scope}
    \fbKGframe{H}
    \draw[quark] (-0.45,0.3) -- (-0.45,1.96);
    \draw[gluon] (0.42,0.33) -- (0.42,1.94);
    \path[mom'={$k$},mom sep=0.32cm] (0.42,0.33) -- (0.42,1.94);
    \node[fb idx,anchor=east] at (0.33,1.9) {$\mu$};
    \node[fb idx,anchor=east] at (0.34,0.5) {$\nu$};
    \node[fb lab] at (0,-0.75) {$g_{\mu\nu}$};
  \end{scope}
  \node[fb op] at (1.85,1.15) {$=$};
  % K part: scalar polarization k_mu at the hard end (dot)
  \begin{scope}[xshift=3.7cm]
    \fbKGframe{H}
    \draw[quark] (-0.45,0.3) -- (-0.45,1.96);
    \draw[longit] (0.42,0.33) -- (0.42,1.94);
    \fill[blue!60!black] (0.42,1.76) circle[radius=1.7pt];
    \path[mom'={$k$},mom sep=0.32cm] (0.42,0.33) -- (0.42,1.94);
    \node[fb idx,anchor=east] at (0.33,1.9) {$\mu$};
    \node[fb idx,anchor=east] at (0.34,0.5) {$\nu$};
    \node[fb lab] at (0,-0.85) {$K_{\mu\nu}=\dfrac{k_\mu\,n_{B\nu}}{n_B\cdot k-\ii0}$};
  \end{scope}
  \node[fb op] at (5.62,1.15) {$+$};
  % G part
  \begin{scope}[xshift=7.55cm]
    \fbKGframe{H}
    \draw[quark] (-0.45,0.3) -- (-0.45,1.96);
    \draw[gluon] (0.42,0.33) -- (0.42,1.94);
    \path[mom'={$k$},mom sep=0.32cm] (0.42,0.33) -- (0.42,1.94);
    \node[fb idx,anchor=east] at (0.33,1.9) {$\mu$};
    \node[fb idx,anchor=east] at (0.34,0.5) {$\nu$};
    \node[fb lab] at (0,-0.75) {$G_{\mu\nu}=g_{\mu\nu}-K_{\mu\nu}$};
    \node[fb lab] at (0,-1.28) {$H\cdot G\cdot C=H\cdot\big(C-\frac{C^+}{k^+}\,k\big)$};
    \node[fb lab,text=red!55!black] at (0,-1.78) {$\ge\lambda$ 压低};
  \end{scope}
  % ======================= row 2: Ward identity for the K part; the exception =======================
  \draw[draw=gray!50!black,line width=1.0pt,-{Stealth[length=6pt,width=5pt]}] (3.35,-1.42) -- (2.3,-2.5);
  \node[fb lab,anchor=east,align=right] at (2.42,-1.86) {对 $H$ 内一切附着求和\\（Ward 恒等式，\cref{fig:ward-identity}）};
  \begin{scope}[xshift=1.6cm,yshift=-5.85cm]
    \node[jet blob,ellipse,minimum width=1.75cm,minimum height=0.66cm,inner sep=0pt] (A) at (0,0) {};
    \node[font=\small] at (-0.42,0) {$A$};
    \node[hard blob,ellipse,minimum width=1.6cm,minimum height=0.66cm,inner sep=0pt] (H) at (0,2.25) {};
    \node[font=\small] at (H) {$H_0$};
    \draw[photon] (-1.25,3.05) -- (H.150);
    \draw[quark] (H.20) -- (1.25,3.0);
    \draw[quark] (-0.45,0.3) -- (H.236);
    \draw[wilson=0.8] (H.236) -- ++(-24:2.45) coordinate (W);
    \node[fb lab,anchor=west] at (W) {沿 $n_B$ 至 $\infty$};
    \coordinate (Y) at ($(H.236)!0.47!(W)$);
    \draw[longit] (0.42,0.33) -- (Y);
    \node[vertex] at (Y) {};
    \path[mom'={$k$},mom sep=0.3cm,mom pos=0.45] (0.42,0.33) -- (Y);
    \node[fb lab,anchor=north west,inner sep=1pt] at ($(W)+(0.05,-0.22)$) {$\dfrac{g\,t^a\,n_B^\nu}{n_B\cdot k-\ii0}$};
  \end{scope}
  % the exception: a single transverse gluon as the active parton (gamma* g -> q qbar)
  \begin{scope}[xshift=8.25cm,yshift=-5.85cm]
    \node[jet blob,ellipse,minimum width=1.75cm,minimum height=0.66cm,inner sep=0pt] (A) at (0,0) {};
    \node[font=\small] at (-0.42,0) {$A$};
    \node[hard blob,ellipse,minimum width=1.6cm,minimum height=0.66cm,inner sep=0pt] (H) at (0,2.25) {};
    \node[font=\small] at (H) {$H$};
    \draw[photon] (-1.25,3.05) -- (H.150);
    \draw[quark] (H.30) -- (1.25,3.0);
    \draw[quark] (1.35,1.75) -- (H.345);
    \draw[gluon] (0.05,0.33) -- (0.05,1.92);
    \path[mom'={$k$},mom sep=0.32cm] (0.05,0.33) -- (0.05,1.92);
    \node[fb idx,anchor=east] at (-0.05,1.15) {$j$};
    \node[fb idx,anchor=west] at (1.25,3.0) {$q$};
    \node[fb idx,anchor=west] at (1.35,1.75) {$\bar q$};
    \node[fb lab,align=center,anchor=north] (EXC) at (0,-0.45) {例外：活动部分子为单个横向胶子\\$k_\mu H^\mu=0$，G 项领头：\\$k^+A^j-k^jA^+\ \to\ G^{+j}$};
    \begin{scope}[on background layer]\node[fb box,fit={(EXC) (A) (H) (-1.3,3.1) (1.55,3.05)}] {};\end{scope}
  \end{scope}
\end{tikzpicture}
\caption{共线胶子的 Grammer--Yennie 分解（Feynman 规范；$k$ 由 $A$ 共线子图流入硬部分，$\mu$ 为硬端、$\nu$ 为共线端的指标）。
$g_{\mu\nu}=K_{\mu\nu}+G_{\mu\nu}$：K 部分（蓝色，圆点标出带标量极化 $k_\mu$ 的硬端）的共线端只取 $A^+=n_B\cdot A$ 分量；对 $H$ 内一切附着求和后，由 Ward 恒等式，K 胶子改为附着在沿 $n_B$ 的 Wilson 线上，乘以不含该胶子的硬部分 $H_0$（左下；$-\ii0$ 对应末态附着，即 DIS 中未来指向的线，初态附着取 $+\ii0$）；
G 部分满足 $G_{\mu\nu}k^\nu=0$，每多一条 G 胶子至少压低 $\lambda$。例外（右下）：作为活动部分子的单个横向胶子（此处为 $\gamma^*g\to q\bar q$），由 Ward 恒等式 $k_\mu H^\mu=0$，其 G 项是领头的，与 $k^+A^j-k^jA^+$ 组成场强 $G^{+j}$。}
\label{fig:method-kg}
\end{figure}
%%% END fig:method-kg


%%% BEGIN fig:method-sudakov-plane
\begin{figure}[htbp]
\centering
\begin{tikzpicture}[fd,
  fb lab/.style={font=\footnotesize,inner sep=1.5pt},
  fb tick/.style={font=\scriptsize},
  fb grid/.style={draw=gray!45,line width=0.4pt,dash pattern=on 1pt off 2pt},
  fb hyp/.style={draw=black!60,line width=0.6pt},
  fb patch/.style={circle,draw=black!70,line width=0.6pt,minimum size=0.62cm,inner sep=0pt,font=\small},
  fb pole/.pic={\draw[line width=1pt,line cap=round] (-0.085,-0.085)--(0.085,0.085) (-0.085,0.085)--(0.085,-0.085);},
  fb contour/.style={draw=blue!60!black,line width=1.0pt},
  fb carrow/.style={postaction={decorate,decoration={markings,
      mark=at position #1 with {\arrow{Stealth[length=5.5pt,width=4.5pt]}}}}}]
  % ======================= (a) the (ln|k^+|, ln|k^-|) plane =======================
  \begin{scope}[x=1.15cm,y=1.15cm]
    % soft--collinear overlap strips (hatched)
    \fill[pattern=north east lines,pattern color=green!45!black] (-2,-2.15) rectangle (0,-1.85);
    \fill[pattern=north east lines,pattern color=green!45!black] (-2.15,-2) rectangle (-1.85,0);
    % axes, ticks, grid
    \draw[lcaxis] (-2.85,-2.85) -- (1.1,-2.85) node[below=1pt,font=\small] {$\ln|k^+|$};
    \draw[lcaxis] (-2.85,-2.85) -- (-2.85,1.1) node[left=1pt,font=\small] {$\ln|k^-|$};
    \foreach \v/\t in {-2/{\lambda^2Q},-1/{\lambda Q},0/{Q}} {
      \draw[line width=0.5pt] (\v,-2.85) -- ++(0,0.07);
      \node[fb tick,below] at (\v,-2.85) {$\t$};
      \draw[line width=0.5pt] (-2.85,\v) -- ++(0.07,0);
      \node[fb tick,left] at (-2.85,\v) {$\t$};
      \draw[fb grid] (\v,-2.78) -- (\v,0.95);
      \draw[fb grid] (-2.78,\v) -- (0.95,\v);
    }
    % lines of constant k^+ k^-
    \draw[fb hyp] (0.75,-0.75) -- (-0.98,0.98) node[left=1pt,fb lab,text=black!70] {$k^+k^-\sim Q^2$};
    \draw[fb hyp,line width=1.0pt] (-2.62,0.62) -- (0.62,-2.62) node[right=1pt,fb lab,text=black!85] {$k^+k^-\sim m^2$};
    % the family k=(z, lambda^2/z, lambda) between C1 and C2
    \draw[red!60!black,line width=1.6pt,dash pattern=on 4pt off 2pt] (-1.72,-0.28) -- (-0.28,-1.72);
    % regions of the off-shell triangle
    \node[fb patch,fill=gray!35] (H) at (0,0) {$H$};
    \node[fb patch,fill=blue!20] (C1) at (0,-2) {$C_1$};
    \node[fb patch,fill=orange!25] (C2) at (-2,0) {$C_2$};
    \node[fb patch,fill=green!25] (S) at (-2,-2) {$S$};
    \draw[black!70,line width=0.6pt,fill=white] (-1,-1) circle[radius=2.3pt];
    \draw[black!60,line width=0.4pt] (-1.06,-1.06) -- (-1.3,-1.4);
    \node[font=\scriptsize,inner sep=1pt,anchor=north] at (-1.36,-1.4) {$(\lambda,\lambda,\lambda)$};
    \node[fb lab,text=red!55!black,anchor=north west,align=left] at (0.02,-0.97) {$k=(z,\lambda^2/z,\lambda)$\\$\lambda^2\ll z\ll1$：见 (b)};
    \node[fb lab,text=green!30!black] at (-1,-2.42) {软—共线重叠};
    \node[font=\small,anchor=south west] at (-3.75,1.15) {(a)};
  \end{scope}
  % ======================= (b) complex k^- plane =======================
  \begin{scope}[xshift=5.55cm,yshift=-1.4cm]
    \fill[red!12] (0,0) circle[radius=0.42];
    \draw[lcaxis] (-2.3,0) -- (3.2,0) node[right=1pt,font=\small] {$\Real k^-$};
    \draw[lcaxis] (0,-1.25) -- (0,1.85) node[left=1pt,font=\small] {$\Imag k^-$};
    % contour: lower half plane on the left, crossing near k_T^2/(4k^+), upper half plane on the right
    \draw[fb contour,fb carrow=0.18,fb carrow=0.86] (-2.2,-1.0) .. controls (-0.6,-1.0) and (0.55,-0.55) .. (1.1,0)
         .. controls (1.6,0.45) and (2.25,0.95) .. (3.05,0.95);
    \pic at (-0.22,0.2) {fb pole};
    \pic at (2.2,-0.2) {fb pole};
    \fill[blue!60!black] (1.1,0) circle[radius=1.5pt];
    \node[fb lab,anchor=south] at (-0.22,0.3) {$D_1$};
    \node[fb lab,anchor=north] at (2.2,-0.3) {$D_g$};
    \draw[gray!50!black,line width=0.4pt] (1.12,0.07) -- (1.28,1.18) node[above,fb lab] {$\Real k^-\approx\dfrac{k_T^2}{4k^+}$};
    \draw[red!50!black,line width=0.4pt] (-0.32,0.28) -- (-0.62,0.56);
    \node[fb lab,text=red!55!black,anchor=south east,align=right] at (-0.58,0.5) {$k^-\approx0$：\\软近似失效};
    \node[fb lab,anchor=north west,align=left] at (-2.2,-1.42) {$D_1$：$k^-\approx-\dfrac{m^2+k_T^2}{2P}+\ii0$\\[2pt]$D_g$：$k^-=\dfrac{k_T^2-\ii0}{2k^+}$};
    \node[font=\small,anchor=south west] at (-2.3,1.55) {(b)};
  \end{scope}
\end{tikzpicture}
\caption{单圈离壳三角图（$p_i^2=-m^2$，$\lambda=m/Q$）的区域与 Collins 减除中的软—共线重叠（\cref{ex:method-sudakov-subtraction}）。
(a)~$(\ln|k^+|,\ln|k^-|)$ 平面（以 $\ln(1/\lambda)$ 为单位，$k_T$ 取 $k_T^2\sim|k^+k^-|$ 的典型值）：区域 $H$、$C_1$、$C_2$、$S$ 位于正方形的四角，其中软区 $S\sim Q(\lambda^2,\lambda^2,\lambda^2)$ 是\cref{def:ir-scalings}的超软标度（$k\sim m^2/Q$）；
粗斜线 $k^+k^-\sim m^2$ 是 $k_T\sim\lambda Q$ 的在壳线 $k_T^2=2k^+k^-$，经过 $C_1$、$C_2$ 与“软中心”$(\lambda,\lambda,\lambda)$；绿色阴影带为软—共线重叠，减除项在重叠区相消的方式由分组表达式决定。
(b)~对 $k=(z,\lambda^2/z,\lambda)$（$\lambda^2\ll z\ll1$，红色虚线），胶子线 $D_g$ 与夸克线 $D_1$ 的 $k^-$ 极点分居实轴两侧、相距 $\sim\lambda^2Q/z$，围道只能在两极点之间穿过实轴；
软近似只在 $\Real k^-\sim k_T^2/(4k^+)$ 附近穿过实轴的围道上成立，在 $k^-\approx0$ 附近（红色区域）失效（\cref{lem:method-i0-compat}）。}
\label{fig:method-sudakov-plane}
\end{figure}
%%% END fig:method-sudakov-plane


%%% BEGIN fig:method-coeff-minus-coll
\begin{figure}[htbp]
\centering
% photons, photon vertices and the cut, common to all graphs (origin: bottom centre of a graph)
\newcommand{\fbCMCphot}{%
  \draw[photon] (-1.22,1.95) -- (-0.88,1.4);
  \draw[photon] (0.88,1.4) -- (1.22,1.95);
  \node[vertex] at (-0.88,1.4) {}; \node[vertex] at (0.88,1.4) {};
  \draw[cut] (0,-0.18) -- (0,2.02);}
% collinear approximation: operator vertices, Wilson line between them, Born struck quark on top;
% #1 = position of the arrow on the Wilson line
\newcommand{\fbCMCapprox}[1]{%
  \draw[quark] (-0.88,0.92) -- (-0.88,1.4);
  \draw[quark=0.78] (-0.88,1.4) -- (0.88,1.4);
  \draw[quark] (0.88,1.4) -- (0.88,0.92);
  \draw[wilson=#1] (-0.88,0.92) -- (0.88,0.92);
  \node[opvertex] at (-0.88,0.92) {}; \node[opvertex] at (0.88,0.92) {};
  \fbCMCphot}
\begin{tikzpicture}[fd,
  fb lab/.style={font=\footnotesize,align=center,anchor=north,inner sep=1pt},
  fb row/.style={font=\small,anchor=east},
  fb frame/.style={rounded corners=3pt,line width=0.6pt}]
  % ======================= row 1: one-loop DIS cut graphs (gamma* q) =======================
  % (i) ISR squared
  \begin{scope}[xshift=0cm]
    \draw[quark] (-0.88,0) -- (-0.88,0.58);
    \draw[quark] (-0.88,0.58) -- (-0.88,1.4);
    \draw[quark=0.78] (-0.88,1.4) -- (0.88,1.4);
    \draw[quark] (0.88,1.4) -- (0.88,0.58);
    \draw[quark] (0.88,0.58) -- (0.88,0);
    \draw[gluon] (-0.88,0.58) -- (0.88,0.58);
    \fbCMCphot
    \node[vertex] at (-0.88,0.58) {}; \node[vertex] at (0.88,0.58) {};
    \node[font=\scriptsize,below,inner sep=1.5pt] at (-0.88,0) {$p$};
    \node[font=\scriptsize,right,inner sep=1.5pt] at (-1.12,1.98) {$q$};
    \node[fb lab] at (0,-0.36) {初态辐射\\的平方};
  \end{scope}
  % (ii) ISR--FSR interference
  \begin{scope}[xshift=2.9cm]
    \draw[quark] (-0.88,0) -- (-0.88,0.58);
    \draw[quark] (-0.88,0.58) -- (-0.88,1.4);
    \draw[quark=0.3] (-0.88,1.4) -- (0.46,1.4);
    \draw[quark] (0.46,1.4) -- (0.88,1.4);
    \draw[quark] (0.88,1.4) -- (0.88,0);
    \draw[gluon] (-0.88,0.58) to[out=0,in=-90] (0.46,1.4);
    \fbCMCphot
    \node[vertex] at (-0.88,0.58) {}; \node[vertex] at (0.46,1.4) {};
    \node[fb lab] at (0,-0.36) {ISR--FSR 干涉\\（+镜像）};
  \end{scope}
  % (iii) FSR squared
  \begin{scope}[xshift=5.8cm]
    \draw[quark] (-0.88,0) -- (-0.88,1.4);
    \draw[quark] (-0.88,1.4) -- (-0.46,1.4);
    \draw[quark=0.27] (-0.46,1.4) -- (0.46,1.4);
    \draw[quark] (0.46,1.4) -- (0.88,1.4);
    \draw[quark] (0.88,1.4) -- (0.88,0);
    \draw[gluon] (-0.46,1.4) arc[start angle=180,end angle=0,radius=0.46];
    \fbCMCphot
    \node[vertex] at (-0.46,1.4) {}; \node[vertex] at (0.46,1.4) {};
    \node[fb lab] at (0,-0.36) {末态辐射\\的平方};
  \end{scope}
  % (iv) virtual vertex correction
  \begin{scope}[xshift=8.7cm]
    \draw[quark] (-0.88,0) -- (-0.88,0.58);
    \draw[quark] (-0.88,0.58) -- (-0.88,1.4);
    \draw[quark] (-0.88,1.4) -- (-0.35,1.4);
    \draw[quark=0.65] (-0.35,1.4) -- (0.88,1.4);
    \draw[quark] (0.88,1.4) -- (0.88,0);
    \draw[gluon] (-0.88,0.58) to[out=0,in=-90] (-0.35,1.4);
    \fbCMCphot
    \node[vertex] at (-0.88,0.58) {}; \node[vertex] at (-0.35,1.4) {};
    \node[fb lab] at (0,-0.36) {顶点虚修正\\（+镜像）};
  \end{scope}
  % ======================= row 2: collinear approximations T_A Gamma = (PDF graph) x Born =======================
  \begin{scope}[yshift=-4.05cm]
    % (i') ladder graph of the PDF operator
    \begin{scope}[xshift=0cm]
      \draw[quark] (-0.88,0) -- (-0.88,0.42);
      \draw[quark] (-0.88,0.42) -- (-0.88,0.92);
      \draw[quark] (0.88,0.92) -- (0.88,0.42);
      \draw[quark] (0.88,0.42) -- (0.88,0);
      \draw[gluon] (-0.88,0.42) -- (0.88,0.42);
      \fbCMCapprox{0.25}
      \node[vertex] at (-0.88,0.42) {}; \node[vertex] at (0.88,0.42) {};
      \node[font=\scriptsize,below,inner sep=1.5pt] at (-0.88,0) {$p$};
      \node[fb lab] at (0,-0.36) {PDF 算符的\\梯形图};
    \end{scope}
    % (ii') real Wilson-line graph
    \begin{scope}[xshift=2.9cm]
      \draw[quark] (-0.88,0) -- (-0.88,0.42);
      \draw[quark] (-0.88,0.42) -- (-0.88,0.92);
      \draw[quark] (0.88,0.92) -- (0.88,0);
      \draw[gluon] (-0.88,0.42) to[out=0,in=-90] (0.46,0.92);
      \fbCMCapprox{0.22}
      \node[vertex] at (-0.88,0.42) {}; \node[vertex] at (0.46,0.92) {};
      \node[fb lab] at (0,-0.36) {Wilson 线实图\\（+镜像）};
    \end{scope}
    % (iii') no target-collinear region
    \begin{scope}[xshift=5.8cm]
      \node[font=\large] at (0,0.9) {$0$};
      \node[fb lab] at (0,-0.36) {无靶共线区};
    \end{scope}
    % (iv') virtual Wilson-line graph
    \begin{scope}[xshift=8.7cm]
      \draw[quark] (-0.88,0) -- (-0.88,0.42);
      \draw[quark] (-0.88,0.42) -- (-0.88,0.92);
      \draw[quark] (0.88,0.92) -- (0.88,0);
      \draw[gluon] (-0.88,0.42) to[out=0,in=-90] (-0.35,0.92);
      \fbCMCapprox{0.72}
      \node[vertex] at (-0.88,0.42) {}; \node[vertex] at (-0.35,0.92) {};
      \node[fb lab] at (0,-0.36) {Wilson 线虚图\\（+镜像）};
    \end{scope}
  \end{scope}
  % ======================= colour frames, row labels, result =======================
  \begin{scope}[on background layer]
    \foreach \x/\c in {0/green,2.9/blue,5.8/gray,8.7/orange}
      \draw[fb frame,draw=\c!55!black,fill=\c!6] (\x-1.4,-5.3) rectangle (\x+1.4,2.2);
  \end{scope}
  \node[fb row] at (-1.5,0.9) {$\Gamma^{(1)}$};
  \node[fb row] at (-1.5,-3.15) {$-\,T_A\Gamma^{(1)}$};
  \foreach \x in {0,2.9,5.8,8.7} \node[font=\large] at (\x,-1.62) {$-$};
  \draw[decorate,decoration={brace,amplitude=5pt},line width=0.6pt] (10.25,2.15) -- (10.25,-5.25);
  \node[anchor=west,align=left,font=\small] at (10.45,-1.55) {各列之和\\[3pt]$=C_H\Gamma^{(1)}$\\[3pt]$=\dfrac{\as}{2\pi}\,C_2^q$\\[1pt]{\footnotesize（$\mu=Q$）}};
\end{tikzpicture}
\caption{单圈 DIS 系数函数 = 割图减去其共线近似（$\gamma^*q$ 道，Feynman 规范；竖直红色虚线为末态割线，左侧为振幅）。
上行为 $\order{\as}$ 割图 $\Gamma^{(1)}$，下行为其靶共线近似 $T_A\Gamma^{(1)}=f^{(1)}_{q/q}\otimes{\rm Born}$：PDF 算符的割图（$\otimes$ 为算符顶点，双线为沿 $n_B$ 的 Wilson 线）乘以领头阶硬部分；同色框中的图一一对应。
初态辐射的平方对应 PDF 算符的梯形图，ISR--FSR 干涉对应胶子连到割线另一侧 Wilson 线的实图，顶点虚修正对应同侧的 Wilson 线虚图（外线自能无标度，未画出）；末态辐射没有靶共线区，不需减除。
各列之差的和即硬区减除项 $C_H\Gamma^{(1)}$，在 $\mu=Q$ 时为 $\frac{\as}{2\pi}C_2^q$；纯维数正规化中下行各图无标度，只留下 $\MSbar$ 抵消项，$C_A\Gamma^{(1)}=-\frac{\as}{2\pi}\frac1\varepsilon P_{qq}$（VF-DIS-6，VF-PDF-3）。}
\label{fig:method-coeff-minus-coll}
\end{figure}
%%% END fig:method-coeff-minus-coll


%%% BEGIN fig:method-glauber-direction
\begin{figure}[htbp]
\centering
% real (k^+,k^-) plane with the Glauber wedge G_delta = {|k^+ k^-| <= delta k_T^2}; units |k_T| = 1, delta = 0.12
\newcommand{\fbGwedge}{%
  \fill[red!9] plot[domain=1.6:0.075,samples=40] (\x,{0.12/\x})
      -- plot[domain=0.075:1.6,samples=40] ({-\x},{0.12/\x})
      -- plot[domain=1.6:0.075,samples=40] ({-\x},{-0.12/\x})
      -- plot[domain=0.075:1.6,samples=40] (\x,{-0.12/\x}) -- cycle;
  \foreach \s/\t in {1/1,-1/1,1/-1,-1/-1}
    \draw[red!55!black,line width=0.5pt] plot[domain=0.075:1.6,samples=40] ({\s*\x},{\t*0.12/\x});
  \draw[lcaxis] (-1.7,0) -- (1.8,0) node[right=0pt,font=\footnotesize] {$k^+$};
  \draw[lcaxis] (0,-1.7) -- (0,1.8) node[above=0pt,font=\footnotesize] {$k^-$};
  \draw[red!50!black,line width=0.4pt] (1.0,-0.06) -- (1.22,-0.45);
  \node[font=\footnotesize,text=red!50!black,anchor=north,inner sep=1pt] at (1.25,-0.45) {$G_\delta$};
  \fill (0.07,0.07) circle[radius=1.6pt];}
\begin{tikzpicture}[fd,
  fb lab/.style={font=\footnotesize,inner sep=1.5pt},
  fb text/.style={font=\footnotesize,align=center,anchor=north,inner sep=1pt},
  fb pole/.pic={\draw[line width=0.9pt,line cap=round] (-0.075,-0.075)--(0.075,0.075) (-0.075,0.075)--(0.075,-0.075);},
  fb def/.style={line width=1.3pt,-{Stealth[length=6pt,width=5pt]}},
  fb contour/.style={draw=blue!60!black,line width=1.0pt},
  fb carrow/.style={postaction={decorate,decoration={markings,
      mark=at position #1 with {\arrow{Stealth[length=5.5pt,width=4.5pt]}}}}}]
  % ======================= (a) one-component deformation stays in G_delta =======================
  \begin{scope}[x=1.15cm,y=1.15cm]
    \fbGwedge
    \draw[fb def,red!65!black] (0.07,0.12) -- (0.07,1.3);
    \node[fb lab,anchor=west,text=red!55!black] at (0.14,1.3) {$\Imag k^-\sim|\mathbf k_T|$};
    \node[font=\small,anchor=south west] at (-1.75,1.6) {(a)};
    \node[fb text] at (0,-1.87) {只形变 $k^-$（$\kappa\lesssim|\mathbf k_T|$）：\\$|k^+k^-|\sim\lambda^3\ll k_T^2$，\\仍在 $G_\delta$ 中};
  \end{scope}
  % ======================= (b) the soft line's own pole limits a one-component deformation =======================
  \begin{scope}[xshift=5.15cm]
    \draw[lcaxis] (-2.1,0) -- (2.1,0) node[right=0pt,font=\footnotesize] {$\Real k^-$};
    \draw[lcaxis] (0,-1.2) -- (0,1.75) node[above=0pt,font=\footnotesize] {$\Imag k^-$};
    \draw[fb contour,fb carrow=0.08,fb carrow=0.92] (-2.05,0) -- (-1.0,0) to[out=0,in=180] (0.55,0.78)
        to[out=0,in=180] (1.75,0) -- (2.05,0);
    \pic at (0.45,-0.16) {fb pole}; \pic at (0.8,-0.16) {fb pole}; \pic at (1.15,-0.16) {fb pole};
    \node[fb lab,anchor=north] at (0.85,-0.25) {$J_A$ 的极点};
    \pic[red!60!black] at (-1.45,0.16) {fb pole};
    \node[fb lab,anchor=south,text=red!55!black] at (-1.45,0.26) {软线极点};
    \draw[{Stealth[length=3.6pt,width=3pt]}-{Stealth[length=3.6pt,width=3pt]},line width=0.5pt] (0.55,0.04) -- (0.55,0.74);
    \node[fb lab,anchor=west] at (0.59,0.4) {$\kappa$};
    \node[font=\small,anchor=south west] at (-2.1,1.38) {(b)};
    \node[fb text] at (0,-2.15) {$k^+<0$：软线极点 $k^-=\frac{k_T^2-\ii0}{2k^+}$\\在上半平面，只形变 $k^-$ 的\\幅度 $\kappa\lesssim k_T^2/|k^+|$};
  \end{scope}
  % ======================= (c) simultaneous deformation along e_G leaves G_delta =======================
  \begin{scope}[xshift=10.35cm,x=1.15cm,y=1.15cm]
    \fbGwedge
    \draw[fb def,green!45!black] (0.03,0.11) -- (-0.92,1.06);
    \node[fb lab,anchor=north east,text=green!30!black] at (-0.5,0.5) {$e_G$};
    \node[font=\small,anchor=south west] at (-1.75,1.6) {(c)};
    \node[fb text] at (0,-1.87) {沿 $e_G=n_B-n_A$ 同时形变：\\$\Imag k^->0$，$\Imag k^+<0$，\\幅度 $\sim|\mathbf k_T|$：离开 $G_\delta$};
  \end{scope}
  % identity below the panels
  \node[font=\footnotesize,anchor=north] at (5.15,-3.82) {所有分母更类时：\ $(l+\ii t\sigma f e_G)^2-m^2=(l^2-m^2+2t^2f^2)+2\ii t\sigma f\,(l^+-l^-)$，\quad $e_G=(-1^+,+1^-,\mathbf 0_T)$};
\end{tikzpicture}
\caption{离开 Glauber 区需要沿 $e_G=n_B-n_A$ 同时形变两个光锥分量（示意；(a)、(c) 以 $|\mathbf k_T|$ 为单位，红色阴影为 Glauber 楔 $G_\delta=\{|k^+k^-|\le\delta k_T^2\}$，复形变量画在同一平面中）。
(a)~对 Glauber 动量 $k^\pm\sim\lambda^2$、$k_T\sim\lambda$ 只形变 $k^-$，即使形变到 $|\mathbf k_T|$ 量级，仍有 $|k^+k^-|\sim\lambda^3\ll k_T^2$，动量留在 $G_\delta$ 中；
(b)~$k^+<0$ 时软线自身的极点位于上半平面，与 $J_A$ 线的极点分居实轴两侧，把只形变 $k^-$ 的幅度限制在 $\kappa\sim k_T^2/|k^+|$（\cref{lem:method-i0-compat}）；
(c)~振幅一侧（$\sigma=+1$）流入 $J_A$ 的软动量沿 $e_G$ 形变（$\Imag k^->0$、$\Imag k^+<0$），幅度 $\sim|\mathbf k_T|$，离开 $G_\delta$，且使所有分母更类时（VF-REF-7；\cref{prop:method-glauber-deformation}）。}
\label{fig:method-glauber-direction}
\end{figure}
%%% END fig:method-glauber-direction


%%% BEGIN fig:disall-seed
\begin{figure}[htbp]
\centering
\begin{tikzpicture}[fd,
  fb lab/.style={font=\footnotesize,inner sep=1.5pt},
  fb op/.style={font=\large},
  fb hregion/.style={hard blob,fill=gray!22,draw=gray!60!black,line width=0.6pt,dash pattern=on 2.5pt off 1.5pt}]
  % ======================= left: one K gluon attached to the hard part =======================
  \begin{scope}
    \coordinate (V) at (-0.4,1.55);       % photon vertex
    \coordinate (G) at (1.0,2.3);         % attachment of the K gluon on the struck quark
    \coordinate (E) at (2.05,2.86);       % outgoing struck quark
    \begin{scope}[on background layer]
      \draw[fb hregion,rotate around={28.2:(0.3,1.92)}] (0.3,1.92) ellipse[x radius=1.12,y radius=0.42];
    \end{scope}
    \node[jet blob,ellipse,minimum width=1.9cm,minimum height=0.7cm,inner sep=0pt] (A) at (0,0) {};
    \node[font=\small] at (-0.45,0) {$A$};
    \draw[photon] (-1.75,2.3) -- (V);
    \draw[quark] (-0.4,0.34) -- (V);
    \draw[quark] (V) -- (G);
    \draw[quark] (G) -- (E);
    \draw[longit] (0.6,0.33) to[out=72,in=-108] (G);
    \fill[blue!60!black] ($(G)+(-0.06,-0.19)$) circle[radius=1.7pt];
    \node[vertex] at (V) {}; \node[vertex] at (G) {};
    \path[mom'={$q$},mom sep=0.24cm,mom length=0.45cm] (-1.75,2.3) -- (V);
    \path[mom={$l$},mom sep=0.24cm,mom length=0.45cm] (-0.4,0.34) -- (V);
    \path[mom={$p'-k$},mom sep=0.3cm,mom length=0.42cm,mom pos=0.42] (V) -- (G);
    \path[mom={$p'$},mom sep=0.24cm,mom length=0.42cm] (G) -- (E);
    \path[mom'={$k$},mom sep=0.26cm,mom length=0.45cm,mom pos=0.45] (0.6,0.33) to[out=72,in=-108] (G);
    \node[fb lab,text=gray!35!black] at ($(G)+(0.42,-0.36)$) {$H$};
    \node[fb lab,anchor=north] at (0.08,-0.42) {$k_\mu H^\mu_a=g\,t^a\,H_0$};
    \node[font=\large,anchor=east] at (-1.75,1.2) {$\displaystyle\sum_{H\text{ 内的附着}}$};
  \end{scope}
  \node[fb op] at (2.75,1.4) {$=$};
  % ======================= middle: H_0 =======================
  \begin{scope}[xshift=4.55cm]
    \coordinate (V) at (-0.25,1.65);
    \draw[photon] (-1.3,2.35) -- (V);
    \draw[quark] (-0.25,0.6) -- (V);
    \draw[quark] (V) -- (1.15,2.4);
    \node[vertex] at (V) {};
    \path[mom={$l+k$},mom sep=0.24cm,mom length=0.45cm] (-0.25,0.6) -- (V);
    \path[mom={$p'$},mom sep=0.24cm,mom length=0.42cm] (V) -- (1.15,2.4);
    \node[fb lab] at (0.6,1.42) {$H_0$};
  \end{scope}
  \node[fb op] at (6.15,1.4) {$\times$};
  % ======================= right: the Wilson-line vertex of the PDF operator =======================
  \begin{scope}[xshift=7.75cm]
    \coordinate (O) at (-0.35,1.65);     % operator vertex psi(0)
    \coordinate (W) at (1.75,2.8);       % towards infinity along n_B
    \coordinate (Y) at ($(O)!0.5!(W)$);
    \node[jet blob,ellipse,minimum width=1.9cm,minimum height=0.7cm,inner sep=0pt] (A) at (0,0) {};
    \node[font=\small] at (-0.45,0) {$A$};
    \draw[quark] (-0.35,0.34) -- (O);
    \draw[wilson=0.78] (O) -- (W);
    \node[opvertex] at (O) {};
    \draw[longit] (0.55,0.33) to[out=70,in=-110] (Y);
    \node[vertex] at (Y) {};
    \path[mom'={$k$},mom sep=0.26cm,mom length=0.45cm,mom pos=0.45] (0.55,0.33) to[out=70,in=-110] (Y);
    \node[fb lab,anchor=east] at ($(O)+(-0.12,0.05)$) {$\psi(0)$};
    \node[fb lab,anchor=west] at (W) {$\infty n_B$};
    \node[fb lab,anchor=north west] at ($(Y)+(0.12,-0.06)$) {$\ii g\,t^a n_B^\alpha$};
    \node[fb lab,anchor=south east] at ($($(O)!0.5!(Y)$)+(-0.02,0.12)$) {$\dfrac{-\ii}{n_B\cdot k-\ii0}$};
    \node[fb lab,anchor=north] at (0.4,-0.42) {$W^{(1)}[\infty n_B,0]$：未来指向};
  \end{scope}
  % summary formula
  \node[font=\footnotesize,anchor=north] at (4.4,-1.15) {$H_\mu\,\dfrac{k^\mu n_B^\alpha}{n_B\cdot k-\ii0}=\Big[\ii g\,t^an_B^\alpha\,\dfrac{-\ii}{n_B\cdot k-\ii0}\Big]H_0=W^{(1)}[\infty n_B,0]\,H_0$};
\end{tikzpicture}
\caption{DIS 中 PDF Wilson 线的种子（VF-WL-8，\cref{prop:method-ward-tree}）：一个来自靶共线子图 $A$ 的 K 胶子（蓝色，流入硬部分的动量为 $k$，圆点标出带标量极化 $k_\mu$ 的硬端）对硬部分 $H$ 的所有附着之和（领头阶即附着在末态被击夸克 $p'$ 上），
经 Ward 恒等式 $k_\mu H^\mu_a=g\,t^aH_0$ 化为去掉该胶子、入射动量为 $l+k$ 的硬部分 $H_0$，乘以夸克场 $\psi(0)$ 处出发、沿 $n_B$ 的未来指向 Wilson 线 $W^{(1)}[\infty n_B,0]$ 的单胶子项（双线箭头为路径方向 $0\to\infty n_B$）。
程函传播子中的 $-\ii0$ 继承自被替代的被击夸克传播子：$(p'-k)^2+\ii0\simeq-2p'^-k^++\ii0$ 的极点位于 $k^+$ 上半平面；Drell--Yan 中同样的附着落在入射反夸克上，给出 $+\ii0$。}
\label{fig:disall-seed}
\end{figure}
%%% END fig:disall-seed


%%% BEGIN fig:disall-ladder-regions
\begin{figure}[htbp]
\centering
\begin{tikzpicture}[fd,
  fb lab/.style={font=\footnotesize,inner sep=1.5pt},
  fb kern/.style={font=\scriptsize,inner sep=1pt,text=gray!25!black},
  fb R2/.style={region,draw=blue!60!black,line width=0.7pt},
  fb R1/.style={region,draw=orange!75!black,line width=0.7pt},
  fb RH/.style={region,draw=gray!55!black,line width=0.7pt}]
  \def\fbX{1.2}
  % ======================= the two-rung ladder (cut graph; amplitude on the left) =======================
  % left rail (amplitude): p -> rung K_1 -> k_1 -> rung K_2 -> k_2 -> photon vertex
  \draw[quark] (-\fbX,0) -- (-\fbX,1.1);
  \draw[quark] (-\fbX,1.1) -- (-\fbX,2.3);
  \draw[quark] (-\fbX,2.3) -- (-\fbX,3.5);
  % right rail (conjugate amplitude)
  \draw[quark] (\fbX,3.5) -- (\fbX,2.3);
  \draw[quark] (\fbX,2.3) -- (\fbX,1.1);
  \draw[quark] (\fbX,1.1) -- (\fbX,0);
  % struck quark and photons (H_0)
  \draw[quark=0.25] (-\fbX,3.5) -- (\fbX,3.5);
  \draw[photon] (-1.62,4.72) -- (-\fbX,3.5);
  \draw[photon] (\fbX,3.5) -- (1.62,4.72);
  % rungs: real gluons crossing the cut
  \foreach \y in {1.1,2.3} {
    \draw[gluon] (-\fbX,\y) -- (\fbX,\y);
    \node[vertex] at (-\fbX,\y) {}; \node[vertex] at (\fbX,\y) {};}
  \node[vertex] at (-\fbX,3.5) {}; \node[vertex] at (\fbX,3.5) {};
  \draw[cut] (0,-0.35) -- (0,4.75);
  % momenta and kernels
  \path[mom'={$k_1$},mom sep=0.25cm,mom length=0.45cm] (-\fbX,1.1) -- (-\fbX,2.3);
  \path[mom'={$k_2$},mom sep=0.25cm,mom length=0.45cm,mom pos=0.42] (-\fbX,2.3) -- (-\fbX,3.5);
  \node[fb lab,anchor=north] at (-\fbX,-0.02) {$p$};
  \node[fb lab,anchor=north] at (\fbX,-0.02) {$p$};
  \node[fb lab,anchor=north,text=gray!30!black] at (0,-0.42) {$p^2=-m^2$};
  \node[fb lab,anchor=south east] at (-1.62,4.72) {$q$};
  \node[fb kern,anchor=south] at (0.6,1.25) {$K_1$};
  \node[fb kern,anchor=south] at (0.6,2.45) {$K_2$};
  \node[fb kern,anchor=south] at (-0.6,3.53) {$\mathsf H_0$};
  % nested hard subgraphs of the three regions
  \draw[fb R2] (-1.55,3.24) rectangle (1.55,3.98);
  \draw[fb R1] (-1.78,2.04) rectangle (1.78,4.14);
  \draw[fb RH] (-2.01,0.84) rectangle (2.01,4.3);
  \node[font=\scriptsize,text=blue!55!black,anchor=south,inner sep=1pt] at (0.72,3.25) {$R_2$};
  \node[font=\scriptsize,text=orange!65!black,anchor=south east,inner sep=1pt] at (1.76,2.06) {$R_1$};
  \node[font=\scriptsize,text=gray!35!black,anchor=south east,inner sep=1pt] at (1.99,0.86) {$R_H$};
  % ======================= the regions and the grouped remainder =======================
  \node[anchor=north west,align=left,font=\footnotesize] (T) at (2.65,4.3) {%
    区域链 $R_2<R_1<R_H$（虚线框为硬子图）：\\[3pt]
    {\color{blue!55!black}$R_2$}：两级横档都与 $p$ 共线；\\[1pt]
    {\color{orange!65!black}$R_1$}：上级横档为硬，下级横档共线；\\[1pt]
    {\color{gray!35!black}$R_H$}：全部为硬——给出两圈系数函数。\\[6pt]
    分组余项：\\[2pt]
    \quad$r=(1-T_H)(1-T_1)(1-T_2)\,\Gamma=\Delta\mathsf H_0\,\Delta K_2\,\Delta K_1$};
  \node[anchor=north west,align=left,font=\footnotesize,text width=6.1cm,draw=red!50!black,line width=0.5pt,
        rounded corners=2pt,fill=red!3,inner sep=4pt] at ($(T.south west)+(0.05,-0.2)$)
    {单个链项 $T_H(1-T_2)\Gamma$ 在 $U_{R_2}$ 上不可积（在 $k_1\parallel\hat p$ 处有无质量共线奇点），因而分组是必要的。};
\end{tikzpicture}
\caption{DIS 两级梯形图（外夸克离壳 $p^2=-m^2$，内线无质量；竖直红色虚线为末态割线，左侧为振幅）的三个区域 $R_2<R_1<R_H$ 与分组余项。
$\Gamma=\mathsf H_0(k_2)\,K_2(k_2,k_1)\,K_1(k_1,p)$，$\mathsf H_0$ 为含两个光子顶点的顶端夸克线，$K_2$、$K_1$ 为上、下两级横档；虚线框依次为 $R_2$、$R_1$、$R_H$ 的硬子图。
分组余项 $r=\Delta\mathsf H_0\Delta K_2\Delta K_1$（$\Delta$ 表示相应因子减去其共线近似）；包含全硬区域 $R_H$ 才能得到两圈系数函数（\cref{ex:disall-ladder}）。}
\label{fig:disall-ladder-regions}
\end{figure}
%%% END fig:disall-ladder-regions


%%% BEGIN fig:disall-heavy
\begin{figure}[htbp]
\centering
\begin{tikzpicture}[fd,
  fb lab/.style={font=\footnotesize,inner sep=1.5pt},
  fb heavy/.style={draw,solid,line width=1.5pt,postaction={decorate,solid,decoration={markings,
      mark=at position #1 with {\arrow[xshift=3.6pt,line width=0.8pt]{Stealth[length=7pt,width=6pt]}}}}},
  fb heavy/.default=0.5]
  % ======================= (a) LO photon--gluon fusion with massive quarks =======================
  \begin{scope}
    \coordinate (B) at (0,0.75); \coordinate (A) at (0,-0.75);
    \draw[photon] (-1.6,0.75) -- (B);
    \draw[gluon] (-1.6,-0.75) -- (A);
    \draw[fb heavy] (B) -- (1.55,0.75);
    \draw[fb heavy] (A) -- (B);
    \draw[fb heavy] (1.55,-0.75) -- (A);
    \node[vertex] at (A) {}; \node[vertex] at (B) {};
    \path[mom={$q$},mom sep=0.26cm] (-1.6,0.75) -- (B);
    \path[mom'={$k$},mom sep=0.3cm] (-1.6,-0.75) -- (A);
    \path[mom={$p_1$},mom sep=0.28cm] (B) -- (1.55,0.75);
    \path[mom'={$p_2$},mom sep=0.3cm] (A) -- (1.55,-0.75);
    \node[fb lab,anchor=east] at (-1.62,0.75) {$\gamma^*$};
    \node[fb lab,anchor=east] at (-1.62,-0.75) {$g$};
    \node[fb lab,anchor=west] at (1.57,0.75) {$Q$};
    \node[fb lab,anchor=west] at (1.57,-0.75) {$\bar Q$};
    \node[fb lab,anchor=west,align=left] at (2.3,0) {$+$\ 交叉图\\\phantom{$+$\ }（$p_1\leftrightarrow p_2$）};
    \node[font=\small,anchor=north] at (0.85,-1.45) {(a)\ 领头阶 $\gamma^*g\to Q\bar Q$（保留 $m$）};
  \end{scope}
  % ======================= (b) ACOT/CWZ subtraction term =======================
  \begin{scope}[xshift=8.15cm]
    \coordinate (B) at (0,0.75);
    \node[jet blob,ellipse,minimum width=1.55cm,minimum height=0.9cm,inner sep=0pt] (F) at (0,-0.75) {};
    \draw[gluon] (-1.75,-0.75) -- (F.west);
    \draw[photon] (-1.6,0.75) -- (B);
    \draw[fb heavy=0.55] (F.north) -- (B);
    \draw[fb heavy] (B) -- (1.55,0.75);
    \draw[fb heavy=0.6] (1.6,-1.25) -- (F.-25);
    \node[vertex] at (B) {};
    \node[font=\footnotesize] at (F) {$f_{Q/g}$};
    \path[mom={$q$},mom sep=0.26cm] (-1.6,0.75) -- (B);
    \path[mom'={$k$},mom sep=0.3cm] (-1.75,-0.75) -- (F.west);
    \node[fb lab,anchor=east] at (-1.62,0.75) {$\gamma^*$};
    \node[fb lab,anchor=east] at (-1.77,-0.75) {$g$};
    \node[fb lab,anchor=west] at (1.57,0.75) {$Q$};
    \node[fb lab,anchor=west] at (1.62,-1.25) {$\bar Q$};
    \node[fb lab,anchor=east] at (-0.06,0.02) {$Q$};
    \node[fb lab,anchor=west,text=gray!30!black] at (0.14,0.38) {$\hat\sigma(\gamma^*Q\to Q)$};
    \node[font=\small,anchor=north] at (0,-1.45) {(b)\ 减除项 $f_{Q/g}\otimes\hat\sigma(\gamma^*Q\to Q)$};
  \end{scope}
  \node[font=\Large] at (5.45,0) {$-$};
\end{tikzpicture}
\caption{含重夸克质量的 DIS（粗线为质量 $m$ 的重夸克 $Q$）。(a)~领头阶的光子—胶子融合，保留完整的质量依赖，两图之和的平方给出领头阶的胶子道系数；
(b)~变味（ACOT/CWZ）方案中的减除项 ${f_{Q/g}\otimes\hat\sigma(\gamma^*Q\to Q)}$：胶子先在重夸克分布 $f_{Q/g}$（蓝色，靶共线）中分裂为 $Q\bar Q$，其中的 $Q$ 再发生领头阶散射 $\gamma^*Q\to Q$。
二者之差在 $m\to0$ 时抵消 (a) 中的共线对数 $\ln(Q^2/m^2)P_{qg}$，给出对 $m/Q$ 一致有效的系数函数（\cref{thm:disall-heavy}，\cref{ex:disall-heavy-lo}）。}
\label{fig:disall-heavy}
\end{figure}
%%% END fig:disall-heavy


%%% BEGIN fig:coll-sia
\begin{figure}[htbp]
\centering
\begin{tikzpicture}[fd,
  fb lab/.style={font=\small},
  fb note/.style={font=\footnotesize,inner sep=1.5pt},
  fb hadron/.style={draw=black!75,line width=1.6pt,line cap=round}]
  \node[hard blob,circle,minimum size=0.74cm,inner sep=0pt] (H)  at (-2.75,0) {};
  \node[hard blob,circle,minimum size=0.74cm,inner sep=0pt] (Hd) at ( 2.75,0) {};
  \node[jet blob,ellipse,minimum width=3.2cm,minimum height=0.85cm,inner sep=0pt] (J1) at (0,1.65) {};
  \node[jetB blob,ellipse,minimum width=3.2cm,minimum height=0.85cm,inner sep=0pt] (J2) at (0,-1.65) {};
  \node[soft blob,circle,minimum size=0.8cm,inner sep=0pt] (S) at (0,0) {};
  % observed hadron h: final-state line of J_1 crossing the cut
  \draw[fb hadron] (-0.42,2.02) to[out=80,in=180] (0,2.62) to[out=0,in=100] (0.42,2.02);
  % photons (gamma* from e+e- annihilation)
  \draw[photon] (-4.3,0) -- (H.west);
  \draw[photon] (Hd.east) -- (4.3,0);
  \path[mom={$q$},mom sep=0.28cm] (-4.3,0) -- (H.west);
  \path[mom={$q$},mom sep=0.28cm] (Hd.east) -- (4.3,0);
  % fermion line: H -> J_1 -> H^dagger -> J_2 -> H
  \draw[quark] (H.62) -- (J1.192);
  \draw[quark] (J1.348) -- (Hd.118);
  \draw[quark] (Hd.242) -- (J2.12);
  \draw[quark] (J2.168) -- (H.298);
  \path[mom={$k$},mom sep=0.28cm,mom length=0.45cm] (H.62) -- (J1.192);
  % soft gluons, on both sides of the cut
  \draw[gluon] (S.135) -- (-0.62,1.27);
  \draw[gluon] (S.45)  -- ( 0.62,1.27);
  \draw[gluon] (S.225) -- (-0.62,-1.27);
  \draw[gluon] (S.315) -- ( 0.62,-1.27);
  % cut
  \draw[cut] (0,3.05) -- (0,-2.45);
  % labels
  \node[fb lab] at (-0.95,1.65) {$J_1$};
  \node[fb lab] at (-0.95,-1.65) {$J_2$};
  \node[fb lab] at (-0.64,0) {$S$};
  \node[fb lab] at (H) {$H$};
  \node[fb lab] at (Hd) {$H^\dagger$};
  \node[fb lab,anchor=south west] at (0.08,2.66) {$h(p_h)$};
  \node[font=\footnotesize,text=red!65!black,anchor=south east] at (-0.05,2.95) {末态割线};
  \draw[blue!50!black,line width=0.4pt] (-1.45,1.98) -- (-2.05,2.45);
  \node[fb note,text=blue!50!black,anchor=south east,align=right] at (-2.0,2.4) {碎裂函数\\$d_{h/q}(z/y)$};
  \node[fb note,anchor=west,align=left] at (1.75,-1.75) {未观测的喷注 $J_2$ 与软辐射 $S$：\\对其末态求和，幺正性相消};
\end{tikzpicture}
\caption{$e^+e^-\to\gamma^*\to hX$ 的领头区域（约化割图；竖直虚线为末态割线，左侧为振幅、右侧为共轭振幅）。
硬子图 $H$、$H^\dagger$ 产生两个喷注：被观测强子 $h$（动量 $p_h$，$z=2p_h\cdot q/q^2$，粗线穿过割线）所在的喷注 $J_1$ 给出碎裂函数 $d_{h/q}(z/y)$，$y$ 为碎裂部分子（动量 $k$）的相应份额；
未观测的喷注 $J_2$ 与软子图 $S$ 也穿过割线，它们在对末态的割线求和中由幺正性相消（\cref{lem:coll-unobserved-jet}），软因子为 $1$；全部共线线都是出射的，Glauber 区没有钳夹（\cref{lem:coll-final-poles}）。}
\label{fig:coll-sia}
\end{figure}
%%% END fig:coll-sia


%%% BEGIN fig:coll-sidis
\begin{figure}[htbp]
\centering
\begin{tikzpicture}[fd,
  fb lab/.style={font=\small},
  fb note/.style={font=\footnotesize,inner sep=1.5pt},
  fb hadron/.style={draw=black!70,line width=1.5pt,line cap=butt},
  fb harc/.style={draw=black!75,line width=1.6pt,line cap=round}]
  \coordinate (HL) at (-1.05,0); \coordinate (HR) at (1.05,0);
  % hard radiation crossing the cut (needed for P_hT ~ Q)
  \draw[gluon] (HL) -- (HR);
  % parton lines: A -> H -> current jet -> H^dagger -> A
  \draw[quark] (-1.05,-1.63) -- (HL);
  \draw[quark] (HL) -- (-1.05,1.63);
  \draw[quark] (1.05,1.63) -- (HR);
  \draw[quark] (HR) -- (1.05,-1.63);
  % photons
  \draw[photon] (-2.45,0.75) -- (HL);
  \draw[photon] (HR) -- (2.45,0.75);
  % target hadron A and the PDF; current jet and the FF
  \draw[fb hadron] (-2.6,-2.0) -- (2.6,-2.0);
  \draw[jet blob] (0,-2.0) ellipse[x radius=1.6,y radius=0.43];
  \draw[jetB blob] (0,2.0) ellipse[x radius=1.6,y radius=0.43];
  \draw[fb harc] (-0.42,2.4) to[out=80,in=180] (0,2.95) to[out=0,in=100] (0.42,2.4);
  \draw[hard blob] (HL) circle[radius=0.36];
  \draw[hard blob] (HR) circle[radius=0.36];
  \draw[cut] (0,-2.75) -- (0,3.3);
  % labels
  \node[fb lab] at (HL) {$H$}; \node[fb lab] at (HR) {$H^\dagger$};
  \node[font=\footnotesize] at (-0.78,-2.0) {$f_{i/A}(\xi)$};
  \node[font=\footnotesize] at (-0.78,2.0) {$d_{h/j}(\zeta)$};
  \node[fb lab,anchor=south west] at (0.07,2.98) {$h(p_h)$};
  \node[fb lab,anchor=east] at (-2.62,-2.0) {$A$}; \node[fb lab,anchor=west] at (2.62,-2.0) {$A$};
  \node[fb lab,anchor=east] at (-2.47,0.75) {$\gamma^*(q)$};
  \path[mom={$\xi P_A$},mom pos=0.42,mom sep=0.26cm,mom length=0.5cm] (-1.05,-1.63) -- (HL);
  \path[mom={$p_h/\zeta$},mom pos=0.6,mom sep=0.26cm,mom length=0.5cm] (HL) -- (-1.05,1.63);
  \node[fb note,anchor=north,inner sep=1.5pt] at (0.36,-0.15) {$X_H$};
  % kinematics and the two regimes
  \node[anchor=west,align=left,font=\footnotesize] (K) at (3.45,1.55) {%
    $x=\dfrac{Q^2}{2P_A\cdot q}\le\xi$，\quad$z=\dfrac{P_A\cdot p_h}{P_A\cdot q}\le\zeta$\\[6pt]
    $P_{hT}$：$h$ 相对 $\gamma^*$--靶轴的横动量};
  \node[anchor=north west,align=left,font=\footnotesize,draw=blue!45!black,fill=blue!4,rounded corners=2pt,
        line width=0.5pt,inner sep=4pt] (C1) at ($(K.south west)+(0.05,-0.35)$)
    {$P_{hT}\sim Q$（共线）：\\$\dd\sigma=\sum_{ij}f_{i/A}\otimes\dd\hat\sigma_{ij}\otimes d_{h/j}$（\cref{thm:coll-sidis}）};
  \node[anchor=north west,align=left,font=\footnotesize,draw=orange!70!black,fill=orange!6,rounded corners=2pt,
        line width=0.5pt,inner sep=4pt] at ($(C1.south west)+(0,-0.25)$)
    {$P_{hT}\ll Q$（TMD，\cref{ch:tmd}）：\\横动量依赖的分布与碎裂函数、软因子};
\end{tikzpicture}
\caption{半单举 DIS $\ell A\to\ell' hX$ 的共线因子化（割图示意；竖直虚线为末态割线，左侧为振幅）。
靶 $A$ 的共线子图给出部分子分布 $f_{i/A}(\xi)$（蓝色，沿 $+z$），流出硬部分的部分子 $j$ 碎裂为被观测强子 $h$（粗线穿过割线），给出碎裂函数 $d_{h/j}(\zeta)$（橙色）；
$P_{hT}\sim Q$ 时 $h$ 的横动量来自硬辐射 $X_H$，截面为 $f\otimes\dd\hat\sigma\otimes d$ 的共线因子化形式；$P_{hT}\ll Q$ 时需要 TMD 因子化，二者在中间区域由 $W+Y$ 衔接（\cref{ch:tmd}）。}
\label{fig:coll-sidis}
\end{figure}
%%% END fig:coll-sidis


%%% BEGIN fig:coll-photoproduction
\begin{figure}[htbp]
\centering
\begin{tikzpicture}[fd,
  fb lab/.style={font=\small},
  fb note/.style={font=\footnotesize,inner sep=1.5pt},
  fb hadron/.style={draw=black!70,line width=1.5pt,line cap=butt},
  fb hb/.style={hard blob,circle,minimum size=0.84cm,inner sep=0pt}]
  % ======================= (a) direct photoproduction =======================
  \begin{scope}
    \node[fb hb] (H) at (0,0) {$H$};
    \node[jet blob,ellipse,minimum width=1.25cm,minimum height=0.66cm,inner sep=0pt] (A) at (-1.55,-1.35) {};
    \node[font=\footnotesize] at (A) {$f_{i/p}$};
    \draw[photon] (-2.75,1.35) -- (H.145);
    \draw[fb hadron] (-3.05,-1.35) -- (A.west);
    \foreach \d in {-0.09,0.09} \draw[plain] ($(A.east)+(0,\d)$) -- ++(3.1,0);
    \draw[quark] (A.45) -- (H.235);
    \draw[quark] (H.25) -- (1.9,0.95);
    \draw[quark] (H.335) -- (1.9,-0.55);
    \node[fb lab,anchor=east] at (-2.77,1.35) {$\gamma$};
    \node[fb lab,anchor=east] at (-3.07,-1.35) {$p$};
    \node[fb note,anchor=west] at (2.2,-1.35) {$X_p$};
    \node[fb note,anchor=west,align=left] at (1.95,0.2) {高 $p_T$\\喷注};
    \path[mom={$\xi P$},mom sep=0.26cm,mom length=0.45cm] (A.45) -- (H.235);
    \node[font=\small,anchor=north,align=center] at (-0.35,-1.85) {(a)\ 直接光子：单强子过程\\{\color{RoyalBlue!80!black}全阶论证 [A]}};
  \end{scope}
  % ======================= (b) resolved photoproduction =======================
  \begin{scope}[xshift=8.0cm]
    \node[fb hb] (H) at (0,0) {$H$};
    \node[jet blob,ellipse,minimum width=1.25cm,minimum height=0.66cm,inner sep=0pt] (A) at (-1.55,-1.35) {};
    \node[jetB blob,ellipse,minimum width=1.25cm,minimum height=0.66cm,inner sep=0pt] (G) at (-1.55,1.35) {};
    \node[font=\footnotesize] at (A) {$f_{i/p}$};
    \node[font=\footnotesize] at (G) {$f_{j/\gamma}$};
    \draw[photon] (-3.05,1.35) -- (G.west);
    \draw[fb hadron] (-3.05,-1.35) -- (A.west);
    \foreach \d in {-0.09,0.09} {
      \draw[plain] ($(A.east)+(0,\d)$) -- ++(3.1,0);
      \draw[plain] ($(G.east)+(0,\d)$) -- ++(3.1,0);}
    \draw[quark] (A.45) -- (H.235);
    \draw[quark] (G.-45) -- (H.125);
    \draw[quark] (H.25) -- (1.75,0.8);
    \draw[quark] (H.335) -- (1.75,-0.8);
    \draw[glauber] (2.0,1.24) -- (2.0,-1.24);
    \node[fb lab,anchor=east] at (-3.07,1.35) {$\gamma$};
    \node[fb lab,anchor=east] at (-3.07,-1.35) {$p$};
    \node[fb note,anchor=west] at (2.2,-1.35) {$X_p$};
    \node[fb note,anchor=west] at (2.2,1.35) {$X_\gamma$};
    \node[fb note,anchor=west,align=left,text=red!55!black] at (2.12,0) {旁观者间的\\Glauber 交换};
    \path[mom={$\xi P$},mom sep=0.26cm,mom length=0.45cm] (A.45) -- (H.235);
    \path[mom'={$x_\gamma q$},mom sep=0.26cm,mom length=0.45cm] (G.-45) -- (H.125);
    \node[font=\small,anchor=north,align=center] at (-0.35,-1.85) {(b)\ 分辨光子：类强子—强子过程\\{\color{Plum!80!black}猜想 [C]}};
  \end{scope}
\end{tikzpicture}
\caption{光生过程 $\gamma p\to\text{jets}+X$（振幅示意）。(a)~直接光子：光子直接参与硬散射 $H$，只有质子一侧有部分子分布 $f_{i/p}$ 与旁观者 $X_p$，与 DIS 同为单强子过程，因子化为全阶论证 [A]（\cref{thm:coll-direct-photo}）；
(b)~分辨光子：光子先“强子化”为由光子部分子分布 $f_{j/\gamma}$ 描述的部分子云，提供进入硬散射的部分子（动量份额 $x_\gamma$），光子一侧也有旁观者 $X_\gamma$，
等效于末态有色的强子—强子过程，旁观者之间可交换 Glauber 胶子（暗红色），因子化为猜想 [C]（\cref{conj:coll-resolved-photo}）。}
\label{fig:coll-photoproduction}
\end{figure}
%%% END fig:coll-photoproduction


%%% BEGIN fig:coll-transversity
\begin{figure}[htbp]
\centering
\begin{tikzpicture}[fd,
  fb lab/.style={font=\small},
  fb note/.style={font=\footnotesize,inner sep=1.5pt},
  fb hadron/.style={draw=black!70,line width=1.5pt,line cap=butt},
  fb harc/.style={draw=black!75,line width=1.6pt,line cap=round},
  fb odd/.style={draw=red!60!black,line width=1.2pt},
  fb oddlab/.style={font=\scriptsize,text=red!55!black,inner sep=1pt}]
  % ======================= (a) inclusive DIS =======================
  \begin{scope}
    \draw[quark] (-0.75,0.3) -- (-0.75,1.55);
    \draw[quark=0.25] (-0.75,1.55) -- (0.75,1.55);
    \draw[quark] (0.75,1.55) -- (0.75,0.3);
    \draw[photon] (-1.45,2.2) -- (-0.75,1.55);
    \draw[photon] (0.75,1.55) -- (1.45,2.2);
    \node[vertex] at (-0.75,1.55) {}; \node[vertex] at (0.75,1.55) {};
    \draw[fb hadron] (-1.9,0) -- (1.9,0);
    \draw[jet blob] (0,0) ellipse[x radius=1.15,y radius=0.4];
    \draw[fb odd] (0,0) ellipse[x radius=1.15,y radius=0.4];
    \node[font=\small] at (-0.52,0) {$h_1$};
    \draw[cut] (0,-0.6) -- (0,2.4);
    \node[fb oddlab,anchor=north] at (0.55,-0.42) {手征奇};
    \node[fb note,anchor=north west,text=gray!30!black] at (-0.72,1.5) {$\gamma^\mu$};
    \node[fb note,anchor=north east,text=gray!30!black] at (0.72,1.5) {$\gamma^\nu$};
    \node[font=\huge,text=red!65!black] at (1.55,0.85) {$\times$};
    \node[fb note,anchor=north,align=center] at (0,-0.85) {(a)\ 单举 DIS：只有一条费米子线，\\硬部分保持手征，$h_1$ 不出现};
  \end{scope}
  % ======================= (b) transversely polarized Drell--Yan =======================
  \begin{scope}[xshift=5.15cm]
    \draw[quark] (-0.75,0.3) -- (-0.75,1.2);
    \draw[quark] (-0.75,1.2) -- (-0.75,2.1);
    \draw[quark] (0.75,2.1) -- (0.75,1.2);
    \draw[quark] (0.75,1.2) -- (0.75,0.3);
    \draw[photon] (-0.75,1.2) -- (-1.65,1.2);
    \draw[photon] (0.75,1.2) -- (1.65,1.2);
    \node[vertex] at (-0.75,1.2) {}; \node[vertex] at (0.75,1.2) {};
    \draw[fb hadron] (-1.9,0) -- (1.9,0);
    \draw[fb hadron] (-1.9,2.4) -- (1.9,2.4);
    \draw[jet blob] (0,0) ellipse[x radius=1.15,y radius=0.4];
    \draw[jetB blob] (0,2.4) ellipse[x radius=1.15,y radius=0.4];
    \draw[fb odd] (0,0) ellipse[x radius=1.15,y radius=0.4];
    \draw[fb odd] (0,2.4) ellipse[x radius=1.15,y radius=0.4];
    \node[font=\small] at (-0.52,0) {$h_1$};
    \node[font=\small] at (-0.52,2.4) {$\bar h_1$};
    \draw[cut] (0,-0.6) -- (0,3.0);
    \node[fb oddlab,anchor=north] at (0.55,-0.42) {手征奇};
    \node[fb oddlab,anchor=south] at (0.55,2.82) {手征奇};
    \node[fb note,anchor=south east] at (-1.67,1.22) {$\gamma^*$};
    \node[font=\Large,text=green!45!black] at (1.62,0.62) {$\checkmark$};
    \node[fb note,anchor=north,align=center] at (0,-0.85) {(b)\ 横向极化 Drell--Yan：\\$h_1\otimes\bar h_1$（两个手征奇因子）};
  \end{scope}
  % ======================= (c) SIDIS with the Collins function =======================
  \begin{scope}[xshift=10.3cm]
    \draw[quark] (-0.75,0.3) -- (-0.75,1.2);
    \draw[quark] (-0.75,1.2) -- (-0.75,2.1);
    \draw[quark] (0.75,2.1) -- (0.75,1.2);
    \draw[quark] (0.75,1.2) -- (0.75,0.3);
    \draw[photon] (-1.65,1.55) -- (-0.75,1.2);
    \draw[photon] (0.75,1.2) -- (1.65,1.55);
    \node[vertex] at (-0.75,1.2) {}; \node[vertex] at (0.75,1.2) {};
    \draw[fb hadron] (-1.9,0) -- (1.9,0);
    \draw[jet blob] (0,0) ellipse[x radius=1.15,y radius=0.4];
    \draw[jetB blob] (0,2.4) ellipse[x radius=1.15,y radius=0.4];
    \draw[fb odd] (0,0) ellipse[x radius=1.15,y radius=0.4];
    \draw[fb odd] (0,2.4) ellipse[x radius=1.15,y radius=0.4];
    \draw[fb harc] (-0.38,2.76) to[out=80,in=180] (0,3.25) to[out=0,in=100] (0.38,2.76);
    \node[font=\small] at (-0.52,0) {$h_1$};
    \node[font=\small] at (-0.52,2.4) {$H_1^\perp$};
    \draw[cut] (0,-0.6) -- (0,3.5);
    \node[fb oddlab,anchor=north] at (0.55,-0.42) {手征奇};
    \node[fb oddlab,anchor=east] at (-1.2,2.4) {手征奇};
    \node[fb note,anchor=south west] at (0.08,3.27) {$h$};
    \node[font=\Large,text=green!45!black] at (1.62,0.62) {$\checkmark$};
    \node[fb note,anchor=north,align=center] at (0,-0.85) {(c)\ SIDIS（Collins 效应，TMD）：\\$h_1\otimes H_1^\perp$};
  \end{scope}
\end{tikzpicture}
\caption{横度分布 $h_1$ 的手征结构（割图示意；深红色边框标出手征奇的分布或碎裂函数）。光子顶点 $\gamma^\mu$ 保持（无质量）夸克的手征，手征奇的因子只能成对出现。
(a)~单举 DIS 只有一条经过硬部分的费米子线，$h_1$ 在领头扭度不出现（翻转手征需要夸克质量，被幂次压低）；
(b)~横向极化 Drell--Yan 中两个手征奇的分布相乘，$h_1\otimes\bar h_1$；(c)~SIDIS 中 $h_1$ 与手征奇的 Collins 碎裂函数 $H_1^\perp$ 相乘（TMD 观测量，见\cref{ch:tmd}）。横度分布的定义见\cref{def:coll-transversity}。}
\label{fig:coll-transversity}
\end{figure}
%%% END fig:coll-transversity
